\documentclass[11pt]{article}
\usepackage[T1]{fontenc}
\usepackage{lmodern}
\usepackage[margin=1in]{geometry}
\usepackage{amsmath,amssymb,amsthm,mathtools}
\usepackage{booktabs,array,microtype}
\usepackage[colorlinks=true,linkcolor=blue,citecolor=blue,urlcolor=blue]{hyperref}
% Public edition used by GitHub and The Kai Project.
\newcommand*{\paperauthor}{Myyura}
\newcommand*{\paperaffiliation}{The Kai Project}
\newcommand*{\paperemail}{}

\hypersetup{pdftitle={Two-colour crossing minimization in complete geometric graphs is NP-hard},pdfauthor={\paperauthor},pdfsubject={NP-completeness and formalized reduction},pdfkeywords={geometric graph, monochromatic crossings, MaxCut, NP-completeness, Lean}}
\newtheorem{theorem}{Theorem}[section]
\newtheorem{lemma}[theorem]{Lemma}
\newtheorem{proposition}[theorem]{Proposition}
\newtheorem{corollary}[theorem]{Corollary}
\theoremstyle{definition}
\newtheorem{definition}[theorem]{Definition}
\theoremstyle{remark}
\newtheorem{remark}[theorem]{Remark}
\numberwithin{equation}{section}
\newcommand{\R}{\mathbb R}
\newcommand{\Q}{\mathbb Q}
\newcommand{\Z}{\mathbb Z}
\newcommand{\N}{\mathbb N}
\newcommand{\1}{\mathbf 1}
\newcommand{\eps}{\varepsilon}
\DeclareMathOperator{\OPT}{OPT}
\DeclareMathOperator{\MaxCut}{MaxCut}
\DeclareMathOperator{\sgn}{sgn}
\DeclareMathOperator{\orient}{or}
\DeclareMathOperator{\dist}{dist}
\DeclarePairedDelimiter{\floor}{\lfloor}{\rfloor}
\DeclarePairedDelimiter{\ceil}{\lceil}{\rceil}
\DeclarePairedDelimiter{\abs}{\lvert}{\rvert}
\DeclarePairedDelimiter{\norm}{\lVert}{\rVert}
\setlength{\emergencystretch}{2em}
\title{Two-colour crossing minimization\protect\\ in complete geometric graphs is NP-hard}
\author{\paperauthor
\ifx\paperaffiliation\empty\else\\{\normalsize\paperaffiliation}\fi
\ifx\paperemail\empty\else\\{\normalsize\href{mailto:\paperemail}{\texttt{\paperemail}}}\fi}
\date{22 October 2026}
\begin{document}
\maketitle
\begin{abstract}
Given a rational point set in general position, colour every edge of its complete straight-line graph with one of two colours and minimize the number of monochromatic properly crossing edge pairs. We prove that the decision problem is NP-complete, even for integer coordinates of polynomial binary length. A deterministic many-one reduction starts from MaxCut on represented permutation graphs. A double-chain comparison grid has two quantitatively stable phases; geometrically realized program points provide controlled fields on a source core. Convex separation excludes grid interactions with program nonfans, while the remaining auxiliary and completion costs are included explicitly. Writing $g(G)=|E(G)|-\MaxCut(G)$, the construction gives a computable baseline $K_{\rm base}$, amplification $\Delta$, and nonnegative errors satisfying
\[
K_{\rm base}+\Delta g(G)-E_-\leq\OPT(P)\leq K_{\rm base}+\Delta g(G)+E_+,
\]
with the strict gap
\[
2^{18}(E_-+E_+)<\Delta.
\]
The new analytic, geometric, and combinatorial parts of the reduction have been formalized in Lean. The formalization explicitly assumes the known convex complete-graph optimum formula; source NP-hardness and the machine-complexity argument are supplied separately. We state the verification boundary and the deterministic coordinate construction in detail.
\end{abstract}
\noindent\textbf{Keywords.} Complete geometric graph; monochromatic crossings; permutation graph; MaxCut; quantitative stability; Lean.
\par\smallskip
\noindent\small\copyright\ 2026 \paperauthor. Licensed under \href{https://creativecommons.org/licenses/by/4.0/}{CC BY 4.0}.\normalsize
\clearpage
\tableofcontents
\clearpage
\section{Introduction}\label{sec:introduction}

Let $P=(p_0,\ldots,p_{N-1})$ be a labelled set of points in the plane.
The complete geometric graph on $P$ contains the straight segment $p_ip_j$
for every $i<j$. We assign one of two colours to every edge and count pairs
of edges that cross properly and have the same colour. Both the point set
and the drawing are fixed; the edge colours are the only variables.

The complexity of this problem for complete geometric graphs was posed at
the Crossing Number Workshop in 2019 and appears as Vogtenhuber's Problem~2
in Section~4.12 of the report of Dagstuhl Seminar~21293
\cite{Dagstuhl2021}. Completeness is a substantial restriction. A reduction
that represents a hard graph by segments does not automatically survive
adding every missing segment between their endpoints: the new edges are
free variables and their optimal contribution can depend on the source
colours. Appendix~\ref{app:completion} gives a six-point example of this
dependence. At the other extreme, for points in convex position the optimum
is the known two-page crossing number of a complete graph
\cite{Abrego2013}.

\begin{definition}\label{def:problem}
The input to \emph{complete geometric two-colour crossing minimization}
is a list of pairwise distinct points $P\in(\Q^2)^N$ with no three
collinear, together with an integer $k$ in the decision version.
Coordinates and $k$ are encoded in binary. Write
\[
 E(P)=\{ij:0\leq i<j<N\},\qquad
 M(P,\chi)=\sum_{\{e,f\}\subseteq E(P)}
       \1[e\text{ and }f\text{ cross properly}]\,\1[\chi(e)=\chi(f)]
\]
for $\chi:E(P)\to\{0,1\}$, and put
$\OPT(P)=\min_\chi M(P,\chi)$.
The decision question is whether $\OPT(P)\leq k$.
\end{definition}

A proper crossing is an intersection of two open segments with four
distinct endpoints. Several crossing pairs may have the same geometric
intersection point; each pair is counted separately. The colouring is
unrestricted: it is not a proper edge colouring in the graph-theoretic
sense. There are no weights or prescribed colours.

\begin{theorem}\label{thm:main}
The decision problem in Definition~\ref{def:problem} is NP-complete under
deterministic polynomial-time many-one reductions. Consequently, computing
$\OPT(P)$ exactly is NP-hard. The decision problem remains NP-complete
when all coordinates are integers encoded in binary.
\end{theorem}

We reduce from \textsc{MaxCut} on permutation graphs, with a permutation
representation supplied as part of the input. Its NP-completeness follows
from de Figueiredo et al.\ \cite{Figueiredo2023}: their construction
explicitly produces the two permutations. After relabelling the first
order, a source graph can be described by a permutation $\pi$ of $[n]$,
with an edge $ij$ when $i<j$ and $\pi(i)>\pi(j)$.
For a source graph $G$, set
\[
 g(G)=\min_{x:V(G)\to\{0,1\}}
          \sum_{uv\in E(G)}\1[x(u)=x(v)]
      =|E(G)|-\MaxCut(G).
\]
Thus a source threshold $k$ becomes the equivalent integer threshold
$g(G)\leq t$, where $t=|E(G)|-k$.

Our construction has three interacting components. A large rectangular
comparison grid forces one of two stable colour phases. Small families of
geometric program points then produce controlled fields on the edges
encoding the source. Finally, an exact convex separation removes the
potentially dominant interaction between the grid and program nonfans.
The remaining added edges are included in an explicit error bound. The
conclusion is a uniform interval
\begin{equation}\label{eq:main-sandwich}
 K_{\mathrm{base}}+\Delta g(G)-E_-
 \ \leq\ \OPT(P)\ \leq\
 K_{\mathrm{base}}+\Delta g(G)+E_+,
 \qquad 2^{18}(E_-+E_+)<\Delta,
\end{equation}
where $K_{\mathrm{base}}$ is computable without an optimization oracle.
The lower bound holds for an unrestricted optimum of the complete graph;
the upper bound is attained within the stated error by a single complete
colouring. No colour restrictions are part of the output instance.

\subsection{Conventions and scales}

We use $N$ for the number of points of the final geometric instance.
The letter $N_g$ will denote the number of vertices of a comparison grid
when a distinction is needed. For nonnegative integers $v$, define
\[
 Z(v)=\frac14\left\lfloor\frac v2\right\rfloor
                  \left\lfloor\frac{v-1}{2}\right\rfloor
                  \left\lfloor\frac{v-2}{2}\right\rfloor
                  \left\lfloor\frac{v-3}{2}\right\rfloor
 \quad(v\geq3),\qquad Z(0)=Z(1)=Z(2)=0.
\]
The theorem of \'{A}brego et al.\ \cite[Theorem~7]{Abrego2013} says that
$Z(v)$ is the optimum for a convex complete geometric graph: a same-page
pair of edges in a two-page drawing crosses exactly when its endpoints
alternate in the linear order. Boundary edges do not affect this count.
We use this established value for two convex background components.

For a source with $n\geq1$, the parameters are
\begin{equation}\label{eq:parameters}
\begin{gathered}
 q=2^{20}(n+1),\quad S=nq,\quad d=S+1,\quad H=2^{20000}d^{40},\\
 W=H^4,\quad B=H^5,\quad D_m=512SW,\quad o=2SW+D_m=514SW,\\
 T=2^{13}o^2,\qquad A_0=2^{13}dT,\qquad
 \Delta=q^2W^4.
\end{gathered}
\end{equation}
All exponents and constants are absolute. These deliberately generous
scales simplify exact error separation; no practical implementation of the
full, extremely large target instances is claimed.

\subsection{Membership in NP and integer coordinates}

\begin{proposition}\label{prop:np}
The problem in Definition~\ref{def:problem} belongs to NP. Rational and
integer coordinate versions are polynomial-time equivalent.
\end{proposition}
\begin{proof}
Let $L$ be the total input bit length. Normalize denominators to be
positive and let $D$ be the product of all coordinate denominators.
Its bit length is $O(L)$. Multiplying every coordinate by this same
positive integer yields integer coordinates of $O(L)$ bits and preserves
all orientation signs. Computing $D$ and the scaled coordinates takes
polynomial time by ordinary integer arithmetic.

Pairwise distinctness and absence of collinear triples can be checked
by examining point pairs and the determinants
\[
 \orient(a,b,c)=(b_x-a_x)(c_y-a_y)-(b_y-a_y)(c_x-a_x).
\]
Declare invalid inputs to be no-instances. A certificate consists of the
$\binom N2$ edge colours. For every unordered pair $ab,cd$ with four
distinct endpoints, a proper crossing is equivalent to
\[
 \orient(a,b,c)\orient(a,b,d)<0,
 \qquad \orient(c,d,a)\orient(c,d,b)<0.
\]
There are $O(N^4)$ such tests, each using a constant number of operations
on integers of polynomial bit length. Count those crossing pairs with
equal colours and compare the result with $k$. Negative thresholds are
immediate no-instances. This verifies a certificate in polynomial time.
The same positive scaling is a reduction from rational to integer inputs;
the converse inclusion is immediate.
\end{proof}

The rest of the paper establishes the reduction and the quantitative
interval~\eqref{eq:main-sandwich}. The formalization and its precise
external dependencies are described in Section~\ref{sec:formalization}.

\section{Stability of the comparison grid}
\label{sec:stability}

The comparison grid supplies a reference colouring without prescribing any
target edge colours.  We first prove a quantitative continuous statement,
then transfer it to every finite rectangular grid.  The finite statement
applies to minimizers with arbitrary bounded unary perturbations.
A separate estimate applies to any colouring with the resulting localized
support; this distinction will allow us to change the program colours while
keeping the same grid colouring.

\subsection{The continuous comparison energy}

All integrals in this section use Lebesgue measure.  Functions are Borel
measurable; the statements also hold for Lebesgue-measurable functions after
replacement on a null set.  For a sign function
$f:[0,1]^2\to\{-1,1\}$, put
\begin{equation}
 \mathcal E(f)=
 \int_{\substack{0\le u<u'\le1\\0\le v<v'\le1}}
 \mathbf 1\{f(u,v)=f(u',v')\}\,du\,dv\,du'\,dv'.
 \label{eq:comparison-energy}
\end{equation}
Each comparable unordered pair occurs once.  Define
\[
 f_*(u,v)=
 \begin{cases}1,&u+v<1,\\-1,&u+v\ge1.\end{cases}
\]
The diagonal convention has no effect on the continuous integral, but we
retain this strict inequality when passing to finite grids.

\begin{theorem}[Continuous two-phase stability]
\label{thm:continuous}
Every measurable sign function on the unit square has
$\mathcal E(f)\ge1/12$.  For every $\varepsilon\ge0$, if
$\mathcal E(f)\le1/12+\varepsilon$, then
\begin{equation}
 \min_{\sigma\in\{-1,1\}}
 \|f-\sigma f_*\|_{L^1([0,1]^2)}
 \le4096\varepsilon^{1/72}.
 \label{eq:continuous-stability}
\end{equation}
\end{theorem}

We give the analytic ingredients with constants, since only a uniform
estimate, rather than a sharp exponent, is needed later.

For $h:[0,1]\to\{-1,1\}$ define
\[
 (\mathcal T h)(t)=\int_0^1\operatorname{sgn}(s-t)h(s)\,ds,
 \qquad D(h)=\frac13-\|\mathcal T h\|_2^2.
\]
Values of the sign kernel on its diagonal are immaterial.  A threshold
function is a sign function that is constant on each side of a point
$c\in[0,1]$, with opposite signs on those two intervals.  Both orientations
and, through $c=0,1$, both constant functions are included.

\begin{lemma}[One-dimensional stability]
\label{lem:one-dimensional}
For every measurable sign function $h$, one has $D(h)\ge0$.  There is a
threshold function $b$ such that
\[
 \|h-b\|_1\le4D(h)^{1/3}.
\]
\end{lemma}

\begin{proof}
Let
\[
 I(h)=\int_{0<r<s<t<1}
 \mathbf 1\{h(r)=h(t)\ne h(s)\}\,dr\,ds\,dt.
\]
The elementary identity
\[
 4\mathbf 1\{h(r)=h(t)\ne h(s)\}
 =1+h(r)h(t)-h(r)h(s)-h(s)h(t)
\]
gives
\begin{equation}
 D(h)=8I(h).
 \label{eq:alternation-deficit}
\end{equation}
To check the coefficient, after integrating out the third variable the
coefficient of $h(u)h(v)$, for $u<v$, in the triple expression is
$2(v-u)-1$.  On the other hand,
\[
 \int_0^1\operatorname{sgn}(u-t)\operatorname{sgn}(v-t)\,dt
 =1-2|u-v|.
\]
The constant triple integral is $1/6$, and the ordered two-variable
integral is half the corresponding symmetric integral.  These observations
give \eqref{eq:alternation-deficit}, including its constant term, by Fubini.
All integrands are bounded, so these changes of order are legitimate.

Let $e$ be the least measure of disagreement with a threshold.  For each
orientation, disagreement is continuous in its threshold position, and
$[0,1]$ is compact, so the minimum is attained.  We claim
$I(h)\ge e^3/64$.  There is nothing to prove if $e=0$.  Otherwise let $p$
be the measure of the positive set.  The constant negative threshold shows
that $p\ge e$.  By continuity of cumulative positive mass choose $a\le b$
so that the positive mass before $a$ and after $b$ is $e/4$ each.

If the negative mass between $a$ and $b$ is at least $e/4$, independently
choosing a positive point before $a$, a negative point between $a,b$, and
a positive point after $b$ gives $I(h)\ge(e/4)^3$.  Otherwise compare with
the threshold positive before $b$ and negative after $b$.  Its disagreement
is at least $e$, whereas the negative mass between $a,b$ plus the positive
mass after $b$ is less than $e/2$.  Thus the negative mass before $a$ is
at least $e/2$.  Comparing with the oppositely oriented threshold at $a$
similarly gives negative mass at least $e/2$ after $b$.  The positive mass
between $a,b$ is $p-e/2\ge e/2$.  Negative--positive--negative triples
therefore give the stronger bound $I(h)\ge e^3/8$.  The claim follows.
By \eqref{eq:alternation-deficit}, $e^3\le8D(h)$; since the sign-function
$L^1$ distance is $2e$, the asserted bound follows.
\end{proof}

Apply $\mathcal T$ in the two coordinates and write
\[
 A(u,v)=\int_0^1\operatorname{sgn}(s-u)f(s,v)\,ds,
 \qquad
 C(u,v)=\int_0^1\operatorname{sgn}(s-v)f(u,s)\,ds,
 \qquad M=\int_{[0,1]^2}f.
\]
Expanding the order indicators and the equality indicator, and using
antisymmetry of the sign kernel, gives
\[
 \mathcal E(f)=\frac{1+M^2-\langle A,C\rangle}{8}.
\]
In particular,
\begin{equation}
 8\bigl(\mathcal E(f)-1/12\bigr)
 =M^2+\frac12\|A-C\|_2^2
 +\frac12\left(\frac13-\|A\|_2^2\right)
 +\frac12\left(\frac13-\|C\|_2^2\right).
 \label{eq:comparison-deficit}
\end{equation}
Lemma~\ref{lem:one-dimensional} on each row and column makes every term on
the right nonnegative.  This proves the lower bound in
Theorem~\ref{thm:continuous}.  For $f_*$ the mean vanishes, both slice
families are thresholds, and direct integration gives $A=C$.  Hence its
energy is exactly $1/12$.

Under an excess bound $\varepsilon$, \eqref{eq:comparison-deficit} gives
\begin{equation}
 |M|\le\sqrt{8\varepsilon},\qquad
 \|A-C\|_2\le4\sqrt\varepsilon,
 \qquad
 \int_0^1 D(f(\cdot,v))\,dv\le16\varepsilon,
 \label{eq:deficit-consequences}
\end{equation}
and the analogous column estimate.  Jensen's inequality and
Lemma~\ref{lem:one-dimensional} show that the average row distance from a
threshold, and also the average column distance, is at most
$4(16\varepsilon)^{1/3}\le11\varepsilon^{1/3}$.  These distance functions are
measurable: one may take the infimum over rational threshold positions in
both orientations, using continuity in the threshold position.  No
measurable selection of optimal thresholds is required.

\begin{lemma}[Approximation by a diagonal profile]
\label{lem:diagonal-approximation}
Suppose $0<h\le1$ and $\mathcal E(f)\le1/12+h^6$.  There is a measurable
$b:[0,2]\to\{-1,1\}$ such that, for $f_1(u,v)=b(u+v)$,
\[
 \|f-f_1\|_1\le128h,
 \qquad \mathcal E(f_1)\le1/12+257h.
\]
\end{lemma}

\begin{proof}
It suffices first to treat $h<1/8$.  Extend $f$ by zero outside the square.
A threshold on $[0,1]$, extended by zero, has total variation at most four.
Approximating each row by a threshold and applying the triangle inequality
shows that translation by $r$ in the first coordinate changes the zero
extension in $L^1$ by at most
$22\varepsilon^{1/3}+4|r|$.  The same holds in the second coordinate.
Here the rowwise inequality can be integrated using the measurable
infimum just described, so no selection issue arises.

Let $f_h$ be the convolution of the zero extension with the uniform
probability density on $[-h,h]^2$.  Integrating the translation bound and
using $\varepsilon=h^6$ yields
\begin{equation}
 \|f-f_h\|_{L^1(\mathbb R^2)}
 \le44h^2+8h\le52h.
 \label{eq:box-error}
\end{equation}
On the interior square $Q_h=[h,1-h]^2$ we also have
\begin{equation}
 \| (\partial_u-\partial_v)f_h\|_{L^1(Q_h)}\le2h.
 \label{eq:directional-bound}
\end{equation}
Indeed, on the open unit square, in the sense of distributions,
\[
 (\partial_u-\partial_v)f
 =\tfrac12\partial_u\partial_v(A-C),
\]
because $\partial_u A=-2f$ and $\partial_v C=-2f$.  The mixed derivative
of the uniform averaging kernel is the signed measure at its four corners
with total variation $h^{-2}$.  Convolution on $Q_h$ therefore gives a bound
$\|A-C\|_1/(2h^2)\le2h$ by \eqref{eq:deficit-consequences}.
Only averaging boxes contained in the original square are used here, so no
distributional boundary term from the zero extension enters this estimate.
The box average of a bounded function is Lipschitz: translating either
coordinate changes the averaging rectangle by area at most a constant
times the translation distance.  Consequently the weak directional
derivative above agrees with its almost-everywhere derivative, and the
one-dimensional fundamental theorem applies on almost every diagonal slice.

For each $s$ average $f_h(u,s-u)$ over the interval of $u$ for which
$(u,s-u)\in Q_h$, and denote the average by $j(s)$.  Put $j(s)=0$ when
that interval is empty.  This defines a measurable function with
$|j|\le1$.  On an interval $I$ of length $\ell$, an absolutely continuous
function $g$ satisfies
\[
 \int_I|g-g_I|
 \le\frac1\ell\int_I\int_I|g(x)-g(y)|\,dx\,dy
 \le\ell\int_I|g'|.
\]
The diagonal coordinates $(s,u)$ have Jacobian one and their intervals
have length at most one.  Integrating this inequality and using
\eqref{eq:directional-bound} gives
$\|f_h-j(u+v)\|_{L^1(Q_h)}\le2h$.  In particular no inverse lower bound
on the length of a short diagonal is required.  Outside $Q_h$ the area is
at most $4h$ and the pointwise difference is at most two.  Together with
\eqref{eq:box-error}, this gives
\[
 \|f-j(u+v)\|_{L^1([0,1]^2)}\le52h+2h+8h=62h.
\]
Round $j$ to a sign, with a fixed convention at zero.  For a sign $z$ and
$r\in[-1,1]$, $|z-\operatorname{sgn}(r)|\le2|z-r|$.  Thus the resulting
$f_1$ has distance at most $124h\le128h$.

For any two sign functions $f,g$, expansion of the equality indicator
gives the useful bound
\begin{equation}
 |\mathcal E(f)-\mathcal E(g)|
 \le\frac12\int_{[0,1]^2}|f-g|
       \bigl(uv+(1-u)(1-v)\bigr)\,du\,dv
 \le\|f-g\|_1.
 \label{eq:energy-lipschitz}
\end{equation}
Hence $\mathcal E(f_1)\le1/12+h^6+128h\le1/12+129h$, which is stronger
than the stated bound.  If $h\ge1/8$, any fixed diagonal sign function
suffices: the distance is at most two, the energy is at most $1/4$, and
both claimed bounds follow.  This completes the lemma for $0<h\le1$.
\end{proof}

\begin{proof}[Proof of Theorem~\ref{thm:continuous}]
Assume first that $\varepsilon>0$, put $h=\varepsilon^{1/6}$, and suppose
$\alpha=257h\le1/4096$.  Lemma~\ref{lem:diagonal-approximation} gives
$f_1(u,v)=b(u+v)$ with energy excess at most $\alpha$.
Let $t=\alpha^{1/4}\le1/8$.
By \eqref{eq:deficit-consequences} there is a row at height
$v\in[t,2t]$ whose one-dimensional deficit is at most $16\alpha/t$.
Its sign profile is $b$ restricted to $[v,v+1]$.  By
Lemma~\ref{lem:one-dimensional}, it has a threshold approximation of
error at most $4(16\alpha/t)^{1/3}<11\alpha^{1/4}$.
Restrict that threshold to $[v,1]$ and extend it to $[0,1]$ as a threshold.
The additional interval has length at most $2t$ and costs at most $4t$.
Thus $b$ on $[0,1]$ is within $15\alpha^{1/4}$ of a threshold $b_0$.
Selecting a row with $v\in[1-2t,1-t]$ in the same way gives a threshold
$b_1$ within $15\alpha^{1/4}$ of $s\mapsto b(1+s)$ on $[0,1]$.

Let $\widetilde b(s)=b_0(s)$ for $s<1$ and
$\widetilde b(s)=b_1(s-1)$ for $s\ge1$, and put
$f_2(u,v)=\widetilde b(u+v)$.  The density of $u+v$ on the square is
$s$ for $0\le s\le1$ and $2-s$ for $1\le s\le2$, so it is at most one.
Consequently
\begin{equation}
 \|f_1-f_2\|_1\le\zeta:=30\alpha^{1/4}.
 \label{eq:two-half-error}
\end{equation}
Each coordinate operator $\mathcal T$ has $L^2$ operator norm at most
one, by Cauchy--Schwarz on each unit interval and Fubini.  Moreover,
$\|f_1-f_2\|_2^2=2\|f_1-f_2\|_1$.  Writing
$\mathcal D=\mathcal T_u-\mathcal T_v$, the triangle inequality followed
by $(x+y)^2\le2x^2+2y^2$ gives
\begin{equation}
 \|\mathcal D f_2\|_2^2
 \le32\alpha+16\zeta\le512\alpha^{1/4}.
 \label{eq:two-half-operator}
\end{equation}

Set $m_0=|\{s\in[0,1]:b_0(s)=b_1(s)\}|$ and
$F(t)=\int_0^t(b_0(s)+b_1(s))\,ds$.  Directly integrating the two
coordinate operators yields
\begin{equation}
 \mathcal D f_2(u,v)=F(v)-F(u),
 \qquad
 \|\mathcal D f_2\|_2^2=2\int_0^1(F-\overline F)^2.
 \label{eq:agreement-variance}
\end{equation}
Because $b_0,b_1$ are thresholds, the nonzero derivative of $F$ is
supported on at most two intervals, and is constantly $2$ or $-2$ on
each.  If $m_0>0$, one of those intervals has length
$\ell\ge m_0/2$.  For any constant $c$, the integral of $(F-c)^2$ on
that interval is at least
\[
 \min_d\int_0^\ell(2s-d)^2\,ds=\ell^3/3.
\]
Taking $c=\overline F$ in \eqref{eq:agreement-variance} and using
\eqref{eq:two-half-operator}, we obtain
$m_0^3\le6144\alpha^{1/4}$ and hence
\begin{equation}
 m_0\le19\alpha^{1/12}.
 \label{eq:agreement-bound}
\end{equation}

Replace $b_1$ by $-b_0$ and call the resulting diagonal function $f_3$.
Since the diagonal density is at most one,
$\|f_2-f_3\|_1\le2m_0$.  If $b_0$ is positive before $c$ and negative
after $c$, then
\[
 \int f_3=\int_0^1(2s-1)b_0(s)\,ds=-2c(1-c).
\]
Reversing the threshold orientation reverses the mean.  The distances
from $f_3$ to the two canonical phases are $2c$ and $2(1-c)$, in some
order, since the density weights of the two halves add to one.
Therefore its distance to a canonical phase is at most
$2|\int f_3|$.  With $d=\|f-f_2\|_1$, the triangle inequality now gives
\[
 \min_\sigma\|f-\sigma f_*\|_1
 \le3d+6m_0+2\left|\int f\right|.
\]
Using $d\le128h+30\alpha^{1/4}$,
\eqref{eq:agreement-bound}, and $|\int f|\le3h^3$, the right-hand side
is at most
\[
 384h+90\alpha^{1/4}+114\alpha^{1/12}+6h^3
 \le594\alpha^{1/12}.
\]
Here $0<h\le1$, $\alpha=257h\le1$, and each smaller power term is
bounded by the displayed coefficient times $\alpha^{1/12}$.
Since $257^{1/12}\le2$, this is at most
$1188h^{1/12}\le4096\varepsilon^{1/72}$.

If $257h>1/4096$, the distance between any two sign functions is at most
two.  The inequality $h^{1/12}\ge1/4$ at the cutoff, and monotonicity
thereafter, give the asserted estimate immediately.  Finally, if
$\varepsilon=0$, apply the positive-excess result for every positive
excess parameter.  The minimum of the distances to the two fixed phases
must then be zero.  One of those two phases attains that minimum, proving
the zero-excess case without a choice of a limiting phase.
\end{proof}

\subsection{Finite rectangular grids with unary costs}

For integers $m,n\ge1$, let $J_{m,n}$ have vertex set
$\{0,\ldots,m-1\}\times\{0,\ldots,n-1\}$, where two vertices are adjacent
exactly when both coordinate orders strictly agree.  In this section write
\[
 V=mn,\qquad \rho=\frac1m+\frac1n,\qquad
 u_i=\frac{i+1/2}{m},\qquad v_j=\frac{j+1/2}{n}.
\]
Thus $V$ is the number of vertices of the comparison graph, not the
number of points of the eventual geometric drawing.  Define
\[
 c^+_{ij}=\mathbf1\{u_i+v_j<1\},\qquad c^-_{ij}=1-c^+_{ij},
 \qquad s^\sigma_{ij}=2c^\sigma_{ij}-1.
\]
In particular a diagonal cell centre has spin $-1$ in the positive phase.

Let $U_x:\{0,1\}\to\mathbb R$ be unary costs, and suppose
\[
 |U_x(1)-U_x(0)|\le\alpha V\quad\hbox{for all }x,
 \qquad \alpha\ge0.
\]
No bound is imposed on their additive constants.  Write
\[
 E(c)=\Phi_J(c)+\sum_xU_x(c_x),
\]
where $\Phi_J$ counts monochromatic edges of $J_{m,n}$.
For a set $F$ of grid vertices, $c\oplus F$ denotes the colouring obtained
by flipping exactly $F$.  Define the exact finite quantities
\begin{equation}
 d_x^\sigma=E(c^\sigma\oplus\{x\})-E(c^\sigma),\qquad
 L_\sigma=E(c^\sigma)+\sum_x\min(0,d_x^\sigma),\qquad
 L=\min(L_+,L_-).
 \label{eq:finite-linear-floor}
\end{equation}
For $w\ge0$ let
\[
 D_w=\{(i,j):|u_i+v_j-1|\le w\},\qquad
 C_w=\{(i,j):u_i+v_j\le w\ \hbox{or}\ u_i+v_j\ge2-w\}.
\]

\begin{theorem}[Finite-grid localization and energy bounds]
\label{thm:grid}
Let $a$ satisfy $\alpha+\rho\le a\le2^{-4096}$.
For every rectangular aspect ratio the following statements hold.
\begin{enumerate}
 \item Every minimizer $c$ of $E$ has a phase $\sigma$ such that, with
 $F=\{x:c_x\ne c_x^\sigma\}$,
 \[
 |F|\le52aV,\qquad F\subseteq D_{220a}\cup C_{220a}.
 \]
 \item The minimum energy belongs to
 \[
 [L-2^{32}a^3V^2,\ L+2^{32}a^3V^2].
 \]
 For either prescribed phase, the explicit set
 $F_\sigma^0=\{x:d_x^\sigma<0\}$ has size at most $52aV$ and
 \[
 E(c^\sigma\oplus F_\sigma^0)\le L_\sigma+2^{32}a^3V^2.
 \]
 \item For every $F\subseteq D_{220a}\cup C_{220a}$,
 \[
 E(c^\sigma\oplus F)\ge L_\sigma-2^{32}a^3V^2.
 \]
 This last conclusion requires no optimality of $c^\sigma\oplus F$.
 In fact it holds for arbitrary real unary costs, without their oscillation
 bound, whenever $\rho\le a\le1/65536$.
\end{enumerate}
For rational unary costs, both $L_\sigma$ and the two prescribed-phase
colourings are computable by polynomially many exact arithmetic operations.
\end{theorem}

\begin{proof}
\emph{From cells to the continuous objective.}
Extend a grid colouring to the sign function $f_c$ constant on its
$1/m$ by $1/n$ cells.  Comparable pairs in distinct rows and columns
contribute exactly $\Phi_J(c)/V^2$.  All other contributions are
nonnegative and their total uncoloured measure is
\[
 \frac14-\frac{\binom m2\binom n2}{m^2n^2}
 =\frac1{4m}+\frac1{4n}-\frac1{4mn}\le\frac\rho4.
\]
We will use the weaker bound
\begin{equation}
 0\le\mathcal E(f_c)-\Phi_J(c)/V^2\le\rho/2.
 \label{eq:cell-comparison}
\end{equation}
The functions $f_{c^+}$ and $f_*$ can disagree only where
$|u+v-1|\le\rho/2$.  This band has area at most $\rho$, whence
$\|f_{c^+}-f_*\|_1\le2\rho$, and the analogous assertion holds for
the negative phase.  By \eqref{eq:energy-lipschitz},
$\mathcal E(f_{c^+})\le1/12+2\rho$.
For a minimizer $c$, comparison with $c^+$ gives
$\Phi_J(c)\le\Phi_J(c^+)+\alpha V^2$.  Therefore
\[
 \mathcal E(f_c)\le1/12+\alpha+\tfrac52\rho\le1/12+5a.
\]
Theorem~\ref{thm:continuous} and the canonical cell discrepancy give a
phase whose disagreement fraction $\mu=|F|/V$ satisfies
\begin{equation}
 \mu\le2048(5a)^{1/72}+\rho\le2^{-20}.
 \label{eq:initial-fraction}
\end{equation}
For the final numerical bound, use $5<8$, $a\le2^{-4096}$, and
$\rho\le a$; the first term is less than $2^{11-4093/72}<2^{-45}$.
This initial estimate follows from the minimum condition.

\emph{Canonical local fields.}
For $p=(u,v)$ let $Q_p$ be the union of the lower-left and upper-right
rectangles comparable to $p$.  With $s=1-u-v$, direct integration gives
\[
 \int_{Q_p}f_*=-s(1-|s|).
\]
For $s\ge0$, the lower-left rectangle contributes $uv$, and the
upper-right rectangle contributes $s^2-(1-u)(1-v)$; reflection gives the
other case.  The union of cells corresponding to the finite neighbours of
a cell centre differs from $Q_p$ only in its row and column strips,
of total area at most $\rho$.  The sign-function discrepancy contributes
at most another $2\rho$.  A single grid flip changes the cost by minus
the vertex spin times its neighbour-spin sum.  Consequently, in either
canonical phase,
\begin{equation}
 \left|\frac{d_x^{\rm grid}}V-\ell(u_i,v_j)\right|\le6\rho,
 \qquad
 \ell(u,v)=|1-u-v|(1-|1-u-v|).
 \label{eq:canonical-field}
\end{equation}
At a diagonal centre the continuous field is zero, so the strict tie
convention causes no change to this bound.  Including the unary cost gives
\begin{equation}
 d_x^\sigma\ge V\ell(u_i,v_j)-6aV.
 \label{eq:field-with-unaries}
\end{equation}
If $t\in[0,1]$ and $t(1-t)\le z$, then either $t\le2z$ or
$1-t\le2z$, by considering $t\le1/2$ and $t\ge1/2$ separately.
Thus a small value of $\ell$ places the vertex in a narrow diagonal strip
or in the two corner regions.

\emph{Finite strip and corner counts.}
A rectangle with side lengths $r,t$ contains at most
\[
 (mr+1)(nt+1)\le V(r+\rho)(t+\rho)
\]
cell centres.  For $\rho\le w\le1$, counting by rows gives
$|D_w|\le(2w+\rho)V\le3wV$, and bounding each corner triangle by
a square gives $|C_w|\le8w^2V$.  Two comparable points of $D_w$, ordered
coordinatewise, have the sum of their coordinate differences at most
$2w$.  For any strip point its strip neighbours therefore lie in two
squares of side $2w$, giving at most $18w^2V$ possibilities.  Hence for
every $F\subseteq D_w\cup C_w$,
\begin{equation}
 \deg_F(x)\le26w^2V\quad(x\in D_w),\qquad
 |E(J_{m,n}[F])|\le83w^3V^2.
 \label{eq:strip-counts}
\end{equation}
For the edge bound, strip--strip, strip--corner, and corner--corner
edges contribute at most $27w^3V^2$, $24w^3V^2$, and $32w^4V^2$,
respectively.  Counting an edge more than once when regions overlap only
increases these upper bounds.  The degree assertion is made only for strip
vertices: corner vertices need not have degree of order $w^2V$.

\emph{Bootstrap from initial stability.}
For a disagreeing vertex $x\in F$, changing its colour back in the
minimizer has cost change
\[
 -d_x^\sigma-2\sum_{y\in F:\,xy\in E(J_{m,n})}s_x^\sigma s_y^\sigma
 \ge0.
\]
Thus $d_x^\sigma\le2\deg_F(x)$, and
\eqref{eq:field-with-unaries} yields
\begin{equation}
 \ell(u_i,v_j)\le6a+2\deg_F(x)/V\qquad(x\in F).
 \label{eq:local-optimality-profile}
\end{equation}
Initially $\deg_F(x)\le\mu V$.  The observation following
\eqref{eq:field-with-unaries} gives
\[
 F\subseteq D_w\cup C_w,
 \qquad w=4\mu+12a.
\]
By \eqref{eq:initial-fraction}, $\rho\le w<1/4$.  For
$x\in F\cap D_w$, \eqref{eq:strip-counts} improves
\eqref{eq:local-optimality-profile} to
$\ell\le6a+52w^2$.  In $D_w$ we have
$\ell\ge|1-u-v|/2$, so
$F\cap D_w\subseteq D_{12a+104w^2}$.  The direct row count and the
corner count now imply
\[
 \mu\le25a+216w^2
 \le25a+6912\mu^2+62208a^2.
\]
Since $\mu\le2^{-20}$ and $a\le2^{-4096}$, we have
$6912\mu^2\le\mu/2$ and $62208a^2\le a$.
Absorbing these terms proves $\mu\le52a$.
Substitution in the definition of $w$ gives $w\le220a$, establishing
the first conclusion.  No small degree bound for corner vertices has
been used in this absorption.

\emph{Exact correction for simultaneous flips.}
For every flip set $F$ and arbitrary real unary costs,
\begin{equation}
 E(c^\sigma\oplus F)-E(c^\sigma)
 =\sum_{x\in F}d_x^\sigma
  +2\sum_{\{x,y\}\in E(J_{m,n}[F])}s_x^\sigma s_y^\sigma.
 \label{eq:simultaneous-flips}
\end{equation}
Indeed, an edge with both endpoints flipped keeps its equality status,
while its two single-flip contributions sum to
$-2s_x^\sigma s_y^\sigma$.  The second sum restores exactly this
discrepancy; unary terms already add independently.

If $F\subseteq D_{220a}\cup C_{220a}$ and
$\rho\le a\le1/65536$, \eqref{eq:strip-counts} bounds the absolute
correction by
\[
 2\cdot83\cdot220^3a^3V^2<2^{32}a^3V^2.
\]
The linear sum is at least $\sum_x\min(0,d_x^\sigma)$.
This proves the third conclusion, including its stronger independence
from the unary oscillation and from optimality.  Applying it to the
localized minimizer proves the lower bound in the second conclusion.

For the upper bound in a prescribed phase, take exactly
$F_\sigma^0=\{x:d_x^\sigma<0\}$.  By
\eqref{eq:field-with-unaries}, this set is contained in
$D_{12a}\cup C_{12a}$, and its cardinality is at most
\[
 (36a+1152a^2)V\le52aV.
\]
Its linear sum in \eqref{eq:simultaneous-flips} is the defining sum for
$L_\sigma$, while its correction is bounded by
$2\cdot83\cdot12^3a^3V^2<2^{32}a^3V^2$.
Choosing the phase with smaller $L_\sigma$ proves the upper bound for
the minimum.  Finally, all quantities in \eqref{eq:finite-linear-floor}
are finite sums and single-vertex comparisons.  They require no search
over colourings and can be evaluated exactly in polynomially many
operations for rational unary costs.
\end{proof}

\subsection{Freezing probes with the grid colouring fixed}

The following lemma is about an arbitrary finite simple graph.  In the
application its vertices are geometric edges and its adjacency relation
is proper crossing.  Thus its conclusion includes interactions between
the probes themselves.

\begin{proposition}[Probe freezing]
\label{prop:freezing}
Let a finite simple graph have vertex partition $X,T,Z$, with induced
subgraph on $X$ equal to $J_{m,n}$.  The objective is the number of
monochromatic graph edges, and all vertices may be coloured freely.
Use $V=mn$ and $\rho=1/m+1/n$.  Suppose every vertex of $X$ has at most
$\alpha V$ neighbours outside $X$, every vertex of $T$ has at most
$\beta V$ neighbours outside $X$, and
\[
 \alpha,\beta\ge0,\qquad \alpha+\rho\le a\le2^{-4096}.
\]
For $p\in T$ and $\sigma\in\{+,-\}$ let
\[
 h_p^\sigma=\sum_{x\in N(p)\cap X}s_x^\sigma.
\]
This is its colour-1 cost minus its colour-0 cost when the grid is
canonical.  Choose $y_p^+=1$ when $h_p^+<0$ and $y_p^+=0$ otherwise,
and put $y_p^-=1-y_p^+$.  These are best responses in both phases,
including their complemented zero-field tie convention.

Put $b=104a+\beta$, and suppose that $Q\ge0$ satisfies
\begin{equation}
 |\{p\in T:|h_p^\sigma|\le bV\}|
 \le Q(b+\rho)V
 \quad\hbox{for both phases}.
 \label{eq:soft-probe-count}
\end{equation}
Then every optimal colouring $\chi$ has a phase $\sigma$ and a colouring
$\psi$ such that only vertices of $T$ are changed, every probe is coloured
$y_p^\sigma$, and
\begin{equation}
 0\le\Phi(\psi)-\Phi(\chi)
 \le E_{\rm fr}:=2Q(b+\rho)bV^2.
 \label{eq:freezing-error}
\end{equation}
The unchanged grid differs from $c^\sigma$ on at most $52aV$ vertices,
all in $D_{220a}\cup C_{220a}$.  If $\OPT_{\rm fr}$ is the minimum over
colourings with the probes in either of the two frozen phases, then
\[
 0\le\OPT_{\rm fr}-\OPT\le E_{\rm fr}.
\]
It suffices for \eqref{eq:soft-probe-count} to have the stronger estimate
$|\{p:|h_p^\sigma|\le tV\}|\le Q(t+\rho)V$ for
$0\le t\le t_0$, together with $b\le t_0$.
\end{proposition}

\begin{proof}
Fix the colours outside $X$ in a global optimum.  Its restriction
to $X$ minimizes the resulting grid objective.  Each conditional unary
oscillation is at most the number of outside neighbours, hence at most
$\alpha V$.  Theorem~\ref{thm:grid} gives a phase $\sigma$ and a
disagreement set $F$ with $|F|\le52aV$ and the stated support.  Keep this
grid colouring fixed for the rest of the proof.

For every probe, its field from $X$ differs from $h_p^\sigma$ by at most
$2|F|\le104aV$.  Under any colouring of all other vertices, its field
from neighbours outside $X$ has absolute value at most $\beta V$.
Thus the full field differs from $h_p^\sigma$ by at most $bV$.
If $|h_p^\sigma|>bV$, the sign of the full field is the same as the
canonical sign, whatever the outside colours are.  Single-vertex
optimality of $\chi$ therefore forces this hard probe already to have
colour $y_p^\sigma$.

Recolour each remaining, soft probe to $y_p^\sigma$, in any order.
At every step the grid remains fixed.  The outside field still has
absolute value at most $\beta V$ even after earlier probe changes, since
all probe--probe neighbours were included in that degree bound.  The
cost increase for the current recolouring is at most the magnitude of
the current full field, which is bounded by
\[
 |h_p^\sigma|+104aV+\beta V\le2bV.
\]
There are at most $Q(b+\rho)V$ soft probes by
\eqref{eq:soft-probe-count}.  Summing the sequential cost changes gives
the upper bound in \eqref{eq:freezing-error}.  The lower bound follows
because $\chi$ is a global optimum and $\psi$ is a feasible colouring.
Only probes were changed, so the same grid and its localization remain.

The frozen feasible set is a subset of the original feasible set and
contains the constructed $\psi$.  Its optimum is therefore between
$\Phi(\chi)$ and $\Phi(\psi)$, proving the final optimum bound.
\end{proof}

The last clause of Theorem~\ref{thm:grid} can now be applied to this
unchanged grid under the newly frozen unary costs.  It applies because
the support is localized, whether or not this grid minimizes the new
conditional objective.  This is the reason for retaining the arbitrary
localized-colouring estimate in addition to the optimum bracket.

\section{Geometric interfaces and effective coordinates}\label{sec:geometry}

We construct three sets of points: originals, frames, and programs. Original--frame edges realize the comparison grid; program--frame edges are the probes to which Proposition~\ref{prop:freezing} will apply. All other edges of the complete graph remain present.

\subsection{Two convex chains}
For rational $D\geq10$, $F\geq1$, and $0<\kappa,\kappa_f\leq1/100$, set
\begin{equation}\label{eq:caps}
 P(y)=(-D-\kappa y^2,y),\qquad A(t)=(F+\kappa_f t^2,t),\qquad -1<y,t<1.
\end{equation}
Choose originals $P_j=P(y_j)$ in increasing parameter order, and frames $A_i=A(t_i)$ in decreasing parameter order, with $0\leq j<o$ and $0\leq i<B$.

\begin{lemma}[The base drawing]\label{lem:double-chain}
The base points are in general position. The edges $A_iP_j$ induce the strict comparison grid:
\[
A_iP_j\sim A_kP_l\quad\Longleftrightarrow\quad (i<k\ \text{and}\ j<l)\ \text{or}\ (k<i\ \text{and}\ l<j).
\]
An edge within either chain does not cross an edge joining the chains.
\end{lemma}
\begin{proof}
For $a<b<c$, direct expansion gives
\begin{align*}
\orient(P(a),P(c),A(t))&=-(c-a)[D+F+\kappa_f t^2+\kappa((a+c)t-ac)]<0,\\
\orient(P(a),P(c),P(b))&=-\kappa(c-a)(b-a)(c-b)<0,\\
\orient(A(a),A(c),P(y))&=(c-a)[D+F+\kappa y^2+\kappa_f((a+c)y-ac)]>0,\\
\orient(A(a),A(c),A(b))&=\kappa_f(c-a)(b-a)(c-b)>0.
\end{align*}
The brackets are at least $D+F-3\kappa$ and $D+F-3\kappa_f$, respectively. These formulas cover every single-chain and mixed triple. With three points on one chain, the middle point lies strictly inside the triangle formed by the other two and any point on the other chain: the displayed chord test and the mixed tests for the two subpairs give its three side tests. This excludes an internal-chain edge crossing a joining edge. Two points on each chain form a convex quadrilateral, since the opposite chain lies strictly on one side of either same-chain chord. Its diagonals give the stated comparison relation, with the two parameter orders reversed.
\end{proof}

There are $S$ early source endpoint clusters and $S$ late clusters, each containing $W$ consecutive originals. Between them lie $D_m=512SW$ dummy originals, so $o=514SW$. Each cluster is an equally spaced interval of parameter diameter $\tau$ around a macro parameter. Early, dummy, and late parameters lie, respectively, in $(-3/4,-5/8)$, $(-1/2,1/2)$, and $(5/8,3/4)$. A matching pairs the early and late clusters; Section~\ref{sec:core} specifies it from the source permutation.

Every original chord has direction $(b-a)(-\kappa(a+b),1)$, while an original--frame vector has horizontal component at least $D+F$ and vertical component of absolute value at most two. These direction cones are separated. Programs will lie in $x<-D/2$, $|y|<2$. The preliminary interfaces below use small displacements and clone diameter $\tau$; the effective construction then chooses both by explicit simultaneous bounds.

\subsection{Ordinary programs}
Take $8\mid B$ and $t_i=1-2(i+1)/(B+1)$. The artificial point $A_*=A(0)$ lies in the strict angular gap between the middle frame rows, as viewed from every original. At a macro point $p=P(a)$, an ordinary program has prototype
\begin{equation}\label{eq:ordinary-prototype}
r=p-\mu_o(A_*-p),\qquad \mu_o>0.
\end{equation}

\begin{lemma}[Ordinary interface]\label{lem:ordinary-table}
For sufficiently small $\mu_o$, followed by sufficiently small $\tau$, every clone $p_s$ of the parent satisfies
\begin{equation}\label{eq:ordinary-own}
rA_k\sim p_sA_i\quad\Longleftrightarrow\quad
\begin{cases}i<k,&k<B/2,\\i>k,&k\geq B/2.\end{cases}
\end{equation}
Outside the parent block, the program fan has the neighborhood of a virtual column inserted at that block. A core edge from a parent clone to a different macro or a dummy crosses exactly the first half of the program fans when its other endpoint is later, and exactly the last half when it is earlier. These tables hold on an open neighborhood of the prototype.
\end{lemma}
\begin{proof}
Put $\epsilon_i=\sgn\orient(r,p_s,A_i)$. The strict central gap gives $\epsilon_i=1$ on the first half and $-1$ on the last half, uniformly for all clones after choosing $\tau$. Both endpoint tests for proper crossing are needed. When $i<k$, they require $\epsilon_i=\epsilon_k=1$; when $i>k$, they require $\epsilon_i=\epsilon_k=-1$. For $i=k$ there is a common endpoint. Here $r$ and $p_s$ lie on the same strict side of each frame chord: on $|y|\leq2$ that chord has supporting coordinate
\[
x=F+\kappa_f((t_i+t_k)y-t_it_k)\geq F-5\kappa_f>0,
\]
whereas both points lie in the left strip. This proves~\eqref{eq:ordinary-own}. In every nonparent column the nonzero macro determinants persist under displacement and cloning, giving the virtual-column prefix or suffix.

For the core table, let $w=P(b)$ with $b\ne a$. Then
$\orient(p,w,r)=-\mu_o\orient(p,w,A_*)$, so $r$ and every frame point lie on opposite sides of $pw$. The other test compares the signs of $\orient(r,A_k,p)$ and $\orient(r,A_k,w)$. The former is negative on the first half and positive on the last; at zero displacement the latter has the sign of $b-a$, and this strict sign persists. This gives the claimed halves for the actual clones. All signs have positive margins, hence hold on an open neighborhood. No half-frame assertion is made for a short edge inside the parent cluster; its exact unary cost is retained later.
\end{proof}

\subsection{Rank programs}
For a selected macro chord $P(a)P(b)$, $a<b$, write $d_0=b-a$ and let $a_s,b_t$, $1\leq s,t\leq W$, be its clone parameters. Choose
\[
0<\mu_r<\frac1{100W},\qquad 0<\tau\leq\frac{\mu_rd_0}{100},\qquad
y_L=a+\mu_rd_0,\quad y_R=b-\mu_rd_0.
\]
The supporting horizontal coordinate of an actual selected chord at height $y$ is
\begin{equation}\label{eq:rank-support}
T_{s,t}(y)=-D-\kappa((a_s+b_t)y-a_sb_t).
\end{equation}

\begin{lemma}[Rank interface]\label{lem:rank-table}
At each internal gap $1\leq\ell<W$ there are open regions for left and right programs. Every left program fan crosses precisely the selected edges with $s>\ell$; every right fan crosses precisely those with $t\leq\ell$. The test is independent of the opposite clone rank.

The grid neighborhood is the exact virtual column after the entire early parent block, or before the entire late parent block, independently of $\ell$. On an incident nonselected core edge to $P(w)$, every fan crosses if $a<w<b$, and none crosses otherwise.
\end{lemma}
\begin{proof}
At the unsteered center $r_0=(1-\theta)P(a)+\theta P(b)$, where $\theta=\mu_r$ or $1-\mu_r$,
\[
\orient(r_0,A,P(a))=\theta\orient(P(a),P(b),A)<0,\quad
\orient(r_0,A,P(b))=(1-\theta)\orient(P(a),A,P(b))>0.
\]
Cloning and sufficiently small horizontal steering preserve these opposite endpoint signs for every frame point. The remaining crossing test is exactly
\begin{equation}\label{eq:rank-cross}
rA\sim P(a_s)P(b_t)\quad\Longleftrightarrow\quad x_r<T_{s,t}(y_r).
\end{equation}
Thus this is an internal segment crossing, not merely a support-line intersection.

At $y_L$, the support increases in $s$ and decreases in $t$. The full ranges over the opposite rank satisfy
\begin{align}
G_{L,\ell}
&=T_{\ell+1,W}(y_L)-T_{\ell,1}(y_L)\notag\\
&=\kappa[(a_{\ell+1}-a_\ell)(b_W-y_L)-(y_L-a_\ell)(b_W-b_1)]
\geq\frac{\kappa\tau d_0}{2W}>0.\label{eq:rank-gaps}
\end{align}
Indeed $a_{\ell+1}-a_\ell\geq\tau/W$, $b_W-y_L\geq0.98d_0$, and $y_L-a_\ell\leq1.005\mu_rd_0<0.011d_0/W$. Thus choose
\[
x_{L,\ell}=\tfrac12(T_{\ell,1}(y_L)+T_{\ell+1,W}(y_L)).
\]
At $y_R$, reflection and rank reversal give
$T_{1,\ell}(y_R)-T_{W,\ell+1}(y_R)\geq\kappa\tau d_0/(2W)$, so
\[
x_{R,\ell}=\tfrac12(T_{W,\ell+1}(y_R)+T_{1,\ell}(y_R))
\]
gives exactly $t\leq\ell$. These prototypes tend to the unsteered centers as $\tau\to0$, retaining the endpoint tests above.

For grid neighbors, marker-to-parent directions remain in the original-chord cone, outside the entire frame cone. The sign $\epsilon_i$ in the preceding proof is now constant over frame rows, positive for an original before the marker and negative for one after it. Consequently
\[
rA_k\sim P_jA_i\quad\Longleftrightarrow\quad
\begin{cases}i<k,&P_j\text{ precedes the marker},\\ i>k,&P_j\text{ follows it}.\end{cases}
\]
The left marker is above its whole parent block and the right marker below its whole parent block. Their displacement remains below the next macro gap, and steering does not change these signs.

For a nonselected incident chord, put $L_{a,w}(y)=-D-\kappa((a+w)y-aw)$. At the left unsteered center,
\[
(r_0)_x-L_{a,w}((r_0)_y)=-\kappa\mu_r(b-a)(b-w).
\]
The support-side test excludes $w>b$. For $w<b$, the other endpoint test excludes $w<a$, since its sign at $P(w)$ tends to the sign of $w-a$, while the near-endpoint sign is negative. Thus precisely $a<w<b$ remains. The right support difference is $\kappa\mu_r(b-a)(a-w)$; its near sign is positive and the other endpoint sign is negative precisely for $w<b$, giving the same interval. Strict macro margins preserve these statements. Parent-cluster short edges and the selected partner are excluded from this last table.
\end{proof}

\begin{lemma}[Nonincident exclusion]\label{lem:nonincident}
A program fan does not cross a core edge whose endpoints both avoid its parent block.
\end{lemma}
\begin{proof}
For endpoints in different nonparent blocks, the fan at the parent macro is noncrossing by Lemma~\ref{lem:double-chain}, with nonzero macro determinants. Small displacement and then cloning preserve the exclusion. For a short edge in one nonparent block, the macro fan has positive distance from that block's center. Choose the displacement and clone neighborhood below this distance. The whole short edge remains disjoint from the fan. This argument avoids applying a strict determinant test to a collapsed zero-length edge.
\end{proof}

\subsection{Canonical response fields}
Let $N_g=Bo$, $\rho=1/B+1/o$, $u_k=(k+1/2)/B$, and $v_j=(j+1/2)/o$. The canonical grid spins are
$\sigma_{kj}=2\1_{u_k+v_j<1}-1$. The field of a fan is its grid cost in colour one minus its grid cost in colour zero. Its canonical response minimizes this cost, with colour zero at a tie. For a virtual column after $r_0$ originals, put $v=r_0/o$; its exact field is
\begin{align}
H_k&=\sum_{j<r_0}\sum_{i<k}\sigma_{ij}+\sum_{j\geq r_0}\sum_{i>k}\sigma_{ij},\notag\\
H_{k+1}-H_k&=\sum_{j<r_0}\sigma_{kj}-\sum_{j\geq r_0}\sigma_{k+1,j}.\label{eq:virtual-field}
\end{align}
There is no contribution from row $k$, whose fan shares a frame endpoint.

\begin{lemma}[Field and soft-row bounds]\label{lem:soft-fields}
Write $s_k=1-u_k-v$. Then
\begin{align}
|H_{k+1}-H_k-o(1-2|s_k|)|&\leq2+2o/B,\notag\\
|H_k/N_g+s_k(1-|s_k|)|&\leq6\rho.\label{eq:virtual-estimates}
\end{align}
A rank program has field $g_k=H_k$. For an ordinary program, choose a virtual position at a parent-block boundary. Then $|g_k-H_k|\leq2BW$, and away from the central ordinary switch the increment error is at most $2W$.

Suppose $B\geq o\geq1000$, $8\mid B$, and $W/o\leq1/514$. For either program type and $0\leq z\leq1/100$,
\begin{equation}\label{eq:soft-rows}
|\{k:|g_k|\leq zBo\}|\leq12Bz+3.
\end{equation}
An early program has response spin $+1$ on the middle $3B/4$ rows; a late program has spin $-1$ there. Each early half-frame sum is at least $B/4$ and its full sum at least $B/2$; the late sums have the opposite bounds. For $s$ programs the total soft count is at most $Q(z+\rho)N_g$, where $Q=12s/o$.
\end{lemma}
\begin{proof}
For $f_u(y)=2\1_{y<1-u}-1$, any midpoint sum over a prefix or suffix of column cells differs from $o$ times the integral by at most one: only the threshold cell matters, including a midpoint tie. Also changing $u$ by $1/B$ changes a suffix integral by at most $2/B$. Since
\[
\int_0^v f_u(y)\,dy-\int_v^1 f_u(y)\,dy=1-2|1-u-v|,
\]
the first estimate follows from~\eqref{eq:virtual-field}. The two rectangles $[0,u_k]\times[0,v]$ and $[u_k,1]\times[v,1]$ differ from the actual cell union only in row $k$, of area $1/B$. The cellwise canonical function differs from $2\1_{x+y<1}-1$ in $L^1$ by at most $2\rho$. Integrating the latter function over the rectangles gives $-s_k(1-|s_k|)$, proving even the bound $3\rho$. Only the $W$ parent columns differ for an ordinary program. Each prefix or suffix has absolute sum at most $B$ and changes by at most one away from its switch. This proves its two error bounds.

For early programs $v\leq1/514$, and for late programs $v\geq1-1/514$. A soft row satisfies
\[
|s_k|(1-|s_k|)\leq z+2W/o+6\rho<0.042<7/64.
\]
It therefore belongs to one of $|s_k|\leq1/8$, $s_k\geq7/8$, or $s_k\leq-7/8$. For early ranks these lie in $u\leq1/8$ or $u\geq3/4$; for late ranks in $u\leq1/4$ or $u\geq7/8$. Each lies away from the ordinary switch. On each interval the main increment has constant sign and magnitude at least $3o/4$. Even the weaker error $6+2W$ leaves magnitude at least $o/2$. A field interval of length $2zBo$ contains at most $4Bz+1$ such rows, proving~\eqref{eq:soft-rows}.

On $u_k\in[1/8,7/8]$, an early rank has $1/9<s_k<8/9$, so $s_k(1-s_k)\geq8/81$. The error $2/514+12/o<0.016$ cannot change the negative field sign, hence the response colour is one. The late case reverses the sign. Each half-frame has $3B/8$ controlled and at most $B/8$ other rows, giving the claimed sums. Finally $12sBz+3s\leq(12s/o)(z+\rho)N_g$.
\end{proof}

A frozen fan of colour $c$ contributes $2c-1$ to the core unary coefficient. Early ordinary halves thus contribute positively, and late halves negatively. Early rank programs add positive suffix steps; late ones add negative prefixes whose disappearance is a positive step. Their full weights lie in $[B/2,B]\subseteq[B/4,B]$. All internal gaps at an endpoint have the same response sequence; copies in one open prototype region retain it, independently of other programs.

\subsection{All remaining edge classes}
Partition the complete graph into $X$ (original--frame), $T_f$ (program--frame), $C$ (original--original), $Y$ (program--original or program--program), and $I$ (frame--frame).
\begin{lemma}[Convex separation]\label{lem:separation}
The programs can all be placed strictly inside $\mathcal H=\{(x,y):\phi(x,y)\leq0\}$, where $\phi(x,y)=x+D+\kappa y^2$. Every open $X$ segment lies outside $\mathcal H$. Consequently $X$ has no crossing with $C$ or $Y$. The frame class $I$ has no crossing-graph neighbor outside itself.
\end{lemma}
\begin{proof}
The set $\mathcal H$ is convex. For $z_\lambda=(1-\lambda)P(a)+\lambda A(t)$,
\begin{equation}\label{eq:cap-ray}
\phi(z_\lambda)=\lambda[D+F+\kappa_ft^2+\kappa(2at-a^2)]+\kappa\lambda^2(t-a)^2.
\end{equation}
Its bracket is positive, so every open grid segment is outside. Substituting $\lambda=-\mu_o$ gives, for $0<\mu_o\leq1$,
$\phi(r)\leq-\mu_o(D+F-7\kappa)<0$ at an ordinary prototype. At a rank height, $a_s<y_r<b_t$ and
\[
\phi(T_{s,t}(y_r),y_r)=-\kappa(y_r-a_s)(b_t-y_r)<0.
\]
Its steered horizontal coordinate is a midpoint of two such supports; $\phi$ is affine in $x$, so the actual prototype is inside. If $\eta_{\rm cap}=\min_r(-\phi(r))>0$, perturbation by at most $\delta\leq\min(1/4,\eta_{\rm cap}/8)$ in both coordinates changes $\phi$ by less than $2\delta$ on $|y|\leq2$. Hence every final program remains inside. All $C,Y$ endpoints lie in $\mathcal H$, so convexity excludes their crossing with $X$.

For the frame component put $\psi(x,y)=F+\kappa_fy^2-x$ and use $\mathcal H_f=\{(x,y):\psi(x,y)\leq0\}$. It contains all $I$ segments. Every nonframe point $r$ has $x_r\leq-D$, $|y_r|\leq2$, and along a frame fan
\[
\psi((1-\lambda)A(t)+\lambda r)
=\lambda[F-x_r+\kappa_f(2ty_r-t^2)]+\kappa_f\lambda^2(y_r-t)^2>0
\]
for $0<\lambda\leq1$, since the bracket is at least $D+F-5\kappa_f$. Thus no fan crosses $I$. The remaining $C,Y$ segments lie in $x\leq-D$, disjoint from $\mathcal H_f\subseteq\{x\geq F\}$.
\end{proof}

The isolated component $I$ has optimum $Z(B)$, not zero. Besides $I$--$I$, the only possible crossing types are
\[
X\!:\!X,\quad X\!:\!T_f,\quad T_f\!:\!T_f,\quad
C\!:\!C,\quad C\!:\!T_f,\quad Y\!:\!T_f,\quad Y\!:\!C,\quad Y\!:\!Y.
\]
All interactions involving multiple programs are retained. Let $G_0$ be the crossing graph after deleting $Y$, and set
\begin{equation}\label{eq:nonfan-budget}
p=Bs,\qquad c=\binom o2,\qquad h=os+\binom s2,\qquad
R_Y=h(p+c)+\binom h2.
\end{equation}
Deletion removes only the last three crossing types, so
\begin{equation}\label{eq:nonfan-delete}
\OPT(G_0)\leq\OPT(P)\leq\OPT(G_0)+R_Y.
\end{equation}
The upper bound extends any colouring of $G_0$ by arbitrary colours on $Y$, charging at most one per deleted constraint.

\subsection{A deterministic rational realization}\label{subsec:realization}
We now give explicit rational scales that instantiate all preceding open-region statements. These are the scales used in the formal development.

Fix $D=10$, $F=1$, and $\kappa=\kappa_f=1/100$. Having fixed the integer counts and source pairing, choose the early macros, physically ordered late macros, dummies, and frames by
\begin{equation}\label{eq:base-parameters}
a_j=-\tfrac34+\frac{j}{8(S+1)},\quad
b_l=\tfrac58+\frac{l}{8(S+1)},\quad
z_v=-\tfrac14+\frac{v}{2(D_m+1)},\quad
t_i=1-\frac{2(i+1)}{B+1}.
\end{equation}
Here $1\leq j,l\leq S$, $1\leq v\leq D_m$, and $0\leq i<B$. Let $\mathcal A$ be the set of early and late macro parameters and all dummy parameters; frame parameters are not included. Define
\begin{equation}\label{eq:explicit-scales}
\eta=\min\bigl(\{1\}\cup\{|a-c|:a,c\in\mathcal A,\ a\ne c\}\bigr),\quad
\mu=\frac{\eta^2}{10000(W+1)},\quad
\tau=\frac{\mu\eta^2}{1000(B+1)}.
\end{equation}
Use $\mu_o=\mu_r=\mu$ and clone parameters
$a_r=a+\tau\vartheta_r$, where $\vartheta_r=(2r-W-1)/(2(W-1))$.
All internal rank prototypes use the midpoint formulas of Lemma~\ref{lem:rank-table}; there are $2S(W-1)$ candidates, including those assigned zero copies later.

\begin{lemma}[Simultaneous quantitative margins]\label{lem:common-scales}
The choices~\eqref{eq:explicit-scales} satisfy all geometric interfaces above for the actual clones and every candidate program. In particular no choice depends on later program multiplicities.
\end{lemma}
\begin{proof}
The set $\mathcal A$ is finite with distinct parameters, so $0<\eta\leq1$. Direct substitution gives
\begin{gather}
0<\mu\leq1/10000,\quad 100W\mu<1,\quad
10000\mu\leq\eta^2,\quad 4800\mu<\eta^2,\notag\\
0<\tau\leq\eta/1000,\quad 1000\tau\leq\mu\eta^2,\quad
3(B+1)\tau/2\leq\mu,\quad
100\tau\leq\mu(b-a)
\label{eq:scale-margins}
\end{gather}
for any selected $a<b$, since $b-a\geq\eta$. Every clone remains in $(-1,1)$ and within $\tau/2$ of its macro. A nonparent original parameter $z$ remains on the same side of the parent macro $a$ as its own macro and obeys $|z-a|\geq\eta/2$. Consequently
\begin{equation}\label{eq:ray-margins}
5\mu\leq|z-a|,\qquad 2\mu(b-a)\leq|z-a|.
\end{equation}
If $z$ also avoids the selected partner macro $b$, then $|b-z|\geq\eta/2$ and
\begin{equation}\label{eq:partner-margin}
20\tau\leq\mu(b-a)|b-z|.
\end{equation}
For two nonparent actual parameters $z,v$, even when they are clones of the same macro, their distances to $a$ are each at least $\eta/2$. Thus
\begin{equation}\label{eq:curvature-margins}
1200\mu<|z-a||v-a|,\qquad
200\mu(b-a)<|z-a||v-a|.
\end{equation}
The second inequality uses $b-a\leq2$.

We check that these numerical inequalities pay for both segment-side tests, rather than only support-line intersections. In the ordinary own-column test, the unperturbed parent determinant has magnitude at least $10\mu|t_i|\geq10\mu/(B+1)$. Moving its original parameter from $a$ to a clone $c$ changes it by at most $25|c-a|$. The bound $3(B+1)|c-a|\leq\mu$ therefore preserves its sign. The constant 25 follows by expanding the determinant difference:
\[
\orient(r,P(c),A(t))-\orient(r,P(a),A(t))
=(c-a)\bigl(r_x-A(t)_x-(a+c)(t-r_y)/100\bigr);
\]
here $|r_x|\leq22$, $|r_y|\leq2$. For a nonparent ray, ordinary displacement changes the determinant by at most $36\mu$, while its macro magnitude is at least $10|z-a|$. Inequality~\eqref{eq:ray-margins} preserves that sign. These bounds establish the full ordinary grid table, including its central switch.

For a core support $L_{c,z}(y)=-10-((c+z)y-cz)/100$, the ordinary offset is exactly
\[
r_x-L_{c,z}(r_y)
=-\frac{(a-c)(a-z)}{100}-\mu C(a,c,z),
\qquad C(a,c,z)=11+\frac{a^2-ac-az}{100},
\]
with $10\leq C(a,c,z)\leq12$ for parameters in $[-1,1]$.
When $c$ is a parent clone, the first term has absolute value at most $\tau/100$; the second is strictly negative and dominates it. Combining with the parent and nonparent ray signs gives the ordinary core half-frame table. For two nonparent endpoints, the product margin~\eqref{eq:curvature-margins} preserves the macro support sign when the parent lies between them; if they are on one side, both ray signs agree. In either case there is no crossing. This also proves nonincidence for a short edge in a different cluster without taking a zero-length limit.

For rank programs, the last inequality of~\eqref{eq:scale-margins} and $100W\mu<1$ are exactly the full-range gap hypotheses in~\eqref{eq:rank-gaps}. The midpoint has horizontal steering from its unsteered macro-chord support at most $\tau/50$, obtained by expanding~\eqref{eq:rank-support} with endpoint deviations at most $\tau/2$. Its $L^\infty$ displacement from its parent is at most $\mu(b-a)$. For either actual selected endpoint $c$, the height gap is at least $\tau$ and
\[
|x_r-P(c)_x|\leq |y_r-c|/10.
\]
Indeed for a support through parameters $u,v$,
\[
L_{u,v}(y)-P(c)_x
=\bigl((c-u)(c-v)-(u+v)(y-c)\bigr)/100;
\]
one endpoint differs from $c$ by at most $\tau$, $|c-u|,|c-v|,|u+v|\leq2$, and averaging preserves the bound. The horizontal distance to every frame point is at least 11, while its vertical distance is at most two. The endpoint determinant thus has the strict sign prescribed by whether $c<y_r$ or $c>y_r$. This proves the endpoint-side test in~\eqref{eq:rank-cross} uniformly over both ranks.

A rank displacement bounded by $\delta_r=\mu(b-a)$ changes a nonparent ray determinant by at most $14\delta_r$, compared with macro magnitude at least $10|z-a|$; use~\eqref{eq:ray-margins}. For a nonselected incident edge from parent clone $c$ to actual $z$, the left support offset has the explicit error bound
\[
\left|r_x-L_{c,z}(r_y)+\frac{\mu(b-a)(b-z)}{100}\right|
\leq\frac{3\tau}{100}.
\]
It is the steering error $\tau/50$ plus clone error at most $\tau/100$. Inequality~\eqref{eq:partner-margin} preserves its sign, so the ray and support tests together give exactly $a<z<b$. The right statement follows by reflection. Finally any rank displacement of size $\delta_r$ changes a nonincident chord offset by at most $2\delta_r$, whereas its initial curvature offset is $|(a-z)(a-v)|/100$. Inequality~\eqref{eq:curvature-margins} again preserves the sign. Endpoints on the same side give equal ray signs, and endpoints straddling the parent put program and frame on the same side of their chord. This proves rank nonincidence in all cases.

The strict inwardness formulas in Lemma~\ref{lem:separation} apply directly to these ordinary and steered rank prototypes. Their heights have absolute value below two. Thus all prototypes are valid simultaneously. The argument uses only fixed base data and the candidate gap index.
\end{proof}

Compute each candidate's exact grid neighbors and response sequence. These determine the ordinary and rank weights. The greedy and ceiling rules of Section~\ref{sec:core} now determine the actual multiplicities and hence $s$. None of the earlier choices is repeated.

Let $\mathcal B$ be the set of all fixed original and frame points, and $\mathcal R$ the finite set of used prototypes. Define positive rational margins
\begin{align*}
\eta_{\rm or}&=\min\bigl(\{1\}\cup
 \{|\orient(u,v,r)|:u,v\in\mathcal B,\ r\in\mathcal R,\ \orient(u,v,r)\ne0\}\bigr),\\
\eta_{\rm cap}&=\min\bigl(\{1\}\cup\{-\phi(r):r\in\mathcal R\}\bigr),\\
\eta_{\rm ht}&=\min\bigl(\{1\}\cup\{2-|r_y|:r\in\mathcal R\}\bigr),\\
L_{\mathcal B}&=1+2\sum_{u\in\mathcal B}(|u_x|+|u_y|),\\
\delta&=\min\{1/4,\eta_{\rm cap}/8,\eta_{\rm ht},\eta_{\rm or}/(8L_{\mathcal B})\}.
\end{align*}
These are the finite-minimum radius formulas used in the formalization. The value one handles empty sets. Every nonzero one-program determinant retains its sign under a coordinate change at most $\delta$: its change is bounded by $\delta(|u_x-v_x|+|u_y-v_y|)<L_{\mathcal B}\delta\leq\eta_{\rm or}/8$. At the insertion points below the displacement is less than $\delta/2$ per coordinate. Hence the cap and height bounds are preserved as well.

An unused zero determinant need not retain its value. Every prescribed crossing has four strict determinant tests. A pair sharing a base endpoint remains noncrossing by definition; each other prescribed exclusion has either a strict same-side support pair or strict same-side fan rays. The quantitative proof above supplies these certificates, including parent-cluster short edges, whose two fan rays have the same nonzero sign. Thus preserving the nonzero determinants suffices.
Let $N=B+o+s$ be the final total point count.
Set $m=N^2+N+2$. For each required copy of prototype $r$, test in lexicographic order the $m^2$ rational candidates
\begin{equation}\label{eq:fixed-grid}
r+\delta\left(\frac{i}{m+1}-\frac12,\frac{j}{m+1}-\frac12\right),\qquad 1\leq i,j\leq m.
\end{equation}
Reject coincidences with existing points and collinearity with any existing pair. At most $\binom N2$ forbidden lines each meet this grid in at most $m$ points, and at most $N$ further points are forbidden. Since $\binom{N}{2}m+N<m^2$, some candidate survives. Induction gives distinct points in general position. The same radius and denominator are used throughout. All required one-program tables and response colours persist; actual fan--fan crossings and short-edge unary costs are evaluated after insertion, while program nonfans are covered by~\eqref{eq:nonfan-budget}.

\begin{proposition}[Effective realization]\label{prop:realization}
The construction is a deterministic polynomial-time algorithm with rational coordinates of polynomial binary length. It preserves all stated interfaces simultaneously and yields general position. A common positive scaling yields integer coordinates of polynomial binary length with the same crossing relation.
\end{proposition}
\begin{proof}
Proposition~\ref{prop:program-count} bounds the actual program count, and $2S(W-1)$ bounds the candidate prototypes. Both, together with $B,o$, are fixed polynomials in the source size. The base coordinates have
$O(\log(B+1)+\log(W+1)+\log(S+1)+\log(D_m+1)+1)$ bits. The fixed formulas for $\eta,\mu,\tau$ use a bounded number of rational operations on base parameters and counts. Prototype coordinates are bounded-degree expressions in these values. Each determinant margin therefore has polynomial bit length; finite minima select existing values. The sum defining $L_{\mathcal B}$ has polynomially many terms, so its denominator bit length is bounded by the sum of their polynomial bit lengths. The same holds for $\delta$. Internal rank gaps use direct formulas, not a chain of nested scales.

The integer greedy recurrence is evaluated with cached cumulative weights. Inserting at most $N$ programs tests $O(N^4)$ candidates per insertion, each against $O(N^2)$ exact equations and determinants, for $O(N^7)$ rational operations. All intermediate bit lengths are polynomial; standard integer arithmetic therefore implements the whole algorithm in polynomial bit time. Earlier inserted points only reject candidates and never redefine later scales. Finally the product of the positive coordinate denominators has bit length at most their sum. Scaling by this product preserves strict determinant signs and gives integer coordinates with polynomial bit length.
\end{proof}

\section{Encoding the source graph in the core}\label{sec:core}

Throughout this section the program fans $T_f$ are fixed to their canonical best responses $t^+$ to the grid. We prove a statement about the entire original--original subsystem with these frozen neighbours. Colours forced here are consequences of minimization in this subsystem; they are never imposed in the target instance.

\subsection{Source matching and amplification}
Let $G$ be a permutation graph on $n$ vertices, with a supplied permutation representation, and write
\[
g(G)=\min_{z:V(G)\to\{0,1\}}\sum_{uv\in E(G)}\1_{z_u=z_v}
=|E(G)|-\MaxCut(G).
\]
This represented source is NP-hard by the construction in \cite{Figueiredo2023}. Place its two endpoint orders in consecutive blocks of the original cap, reversing one order. The resulting chords alternate exactly for adjacent source vertices. Replace each source chord by $q$ macro copies, ordered forward at the early endpoint and backward at the late endpoint. Copies of a vertex are nested; two distinct source classes have either all $q^2$ crossings or none. Write $S=nq$ and label the selected macro chords $p_jq_j$, $1\leq j\leq S$, by early endpoint order.

Each macro endpoint is replaced by $W$ consecutive originals, denoted $p_{j,r},q_{j,t}$, $1\leq r,t\leq W$. There are $D_m=512SW$ dummy points between the endpoint blocks. All original points are in convex position on the same cap, so their chords cross precisely when their endpoints alternate.

Partition the core $C$ into selected edges
\[
\mathcal S=\bigcup_{j=1}^S\mathcal S_j,\qquad
\mathcal S_j=\{p_{j,r}q_{j,t}:1\leq r,t\leq W\},
\]
short edges $\mathcal U$ within one endpoint cluster, dummy--dummy edges $\mathcal D$, and the remaining palette edges $\mathcal F$. Define the proposed palette by
\begin{equation}\label{eq:palette}
f^*(e)=
\begin{cases}
0,&e\text{ joins two early clusters or an early cluster to a dummy},\\
1,&e\text{ joins two late clusters or a dummy to a late cluster},\\
\1_{j<k},&e=p_{j,r}q_{k,t},\ j\ne k.
\end{cases}
\end{equation}
Every palette edge falls in exactly one case.

\begin{lemma}[Independent source amplification]\label{lem:source-amplification}
After deleting the crossing constraints internal to each $\mathcal S_j$, the graph induced on $\mathcal S$ is the independent $qW^2$-fold blowup of $G$. Its minimum monochromatic count is $\Delta g(G)$, where $\Delta=q^2W^4$.
\end{lemma}
\begin{proof}
Consecutive endpoint clusters preserve all crossings between different macro chords. The $q$ macro copies of one vertex remain pairwise noncrossing, and the deletion makes each individual bundle independent. Thus each source vertex becomes an independent class of $qW^2$ variables with identical neighbours; source adjacency gives a complete join.

With colours outside one class fixed, its cost is affine in the number of colour-one variables in that class. One uniform colour is therefore no worse. Perform this operation successively on all classes. The final class-uniform colouring costs $(qW^2)^2$ times a colouring cost of $G$. Conversely every source colouring lifts to that cost. The deleted internal constraints remain in the actual drawing and are restored below.
\end{proof}

\subsection{Unary costs from frozen fans}
Every unary coefficient means cost in colour one minus cost in colour zero; a frozen neighbour of colour $a$ contributes $2a-1$. With $e\sim f$ denoting proper crossing, define
\begin{align}
E_C(x;t^+)
&=\sum_{\substack{\{e,f\}\subseteq C\\e\sim f}}\1_{x_e=x_f}
+\sum_{\substack{e\in C,\ f\in T_f\\e\sim f}}\1_{x_e=t_f^+},
&
M_C&=\min_{x:C\to\{0,1\}}E_C(x;t^+).
\label{eq:core-energy}
\end{align}
The first sum ranges over unordered pairs of distinct edges.

By Lemmas~\ref{lem:ordinary-table}--\ref{lem:soft-fields}, an ordinary program at an early endpoint contributes a positive half-frame weight in $[B/4,B]$ on an incident nonshort core edge; at a late endpoint the weight is negative with magnitude in that interval. Which half is used depends only on endpoint order, not clone ranks. An early rank program at gap $\ell$ adds a positive weight $w_\ell\in[B/4,B]$ on the suffix $r>\ell$. A late program adds $-w_\ell$ on the prefix $t\leq\ell$. No program fan meets a nonincident core edge. Its contribution on any incident edge is bounded in magnitude by $B$. The actual unary on a short edge is computed directly. Response weights depend only on one prototype and the base grid; their computation is independent of other multiplicities and core colours.

\begin{lemma}[Affine palette field]\label{lem:affine-palette}
The palette contribution to a selected edge is
\begin{equation}\label{eq:affine-palette}
\sum_{\substack{f\in\mathcal F\\f\sim p_{j,r}q_{j,t}}}(2f^*(f)-1)
=\beta_j+\gamma_{j,L}(r-1)+\gamma_{j,R}(t-1),
\end{equation}
where the integer coefficients are computable by crossing counts and obey
$|\beta_j|\leq o^2$ and $-o\leq\gamma_{j,L},\gamma_{j,R}<0$.
\end{lemma}
\begin{proof}
A palette edge avoiding both parent clusters contributes independently of $r,t$. One with an endpoint in the early parent crosses according to whether that endpoint precedes or follows $r$. Its other endpoint avoids the late parent, and its colour is constant over the early cluster; its sum is affine in $r$ alone. The late cluster similarly gives a function of $t$. An edge between the two parent clusters is selected, and one inside a cluster is short, so neither contributes a mixed term to the palette field.

The dummy neighbours give exactly
\[
-D_m(r-1)+D_m(W-t).
\]
The early--dummy edges have colour zero and cross when their parent rank is below $r$; the dummy--late edges have colour one and cross when their parent rank is above $t$. Changing either rank changes at most one contribution for each nondummy endpoint, of which there are at most $2SW$. Thus either slope lies in
$[-D_m-2SW,-D_m+2SW]\subseteq[-o,0)$. Finally the number of possible palette neighbours is less than $o^2$, proving the intercept bound.
\end{proof}

\begin{lemma}[Greedy rank correction]\label{lem:rank-correction}
Each slope in Lemma~\ref{lem:affine-palette} can be corrected with at most $4oW/B+4$ rank programs at that endpoint. After correction the field on bundle $j$ is
\begin{equation}\label{eq:corrected-field}
b_j+\eps_{j,L}(r)+\eps_{j,R}(t),\qquad
0\leq\eps_{j,L}(r),\eps_{j,R}(t)<B,\qquad |b_j|\leq4o^2.
\end{equation}
\end{lemma}
\begin{proof}
For $\gamma<0$, let $f(r)=-\gamma(r-1)$, so $f(1)=0$ and $f(W)\leq oW$. Starting with $R(1)=0$, at gap $r-1$ add the least nonnegative number of its weight-$w_{r-1}$ programs that makes the new cumulative weight $R(r)$ at least $f(r)$. Inductively
\begin{equation}\label{eq:greedy}
f(r)\leq R(r)<f(r)+B.
\end{equation}
If copies are added, the final overshoot is smaller than the last weight, at most $B$. If none is added, the preceding overshoot remains below $B$ because the target increased. Each copy weighs at least $B/4$, while $R(W)<oW+B$, giving the stated count. Unoccupied gaps add no copies; there is no extra term proportional to $W$.

At the early endpoint the coefficient is $R(r)$, which cancels the slope with error in $[0,B)$. At the late endpoint the negative prefixes give $-R(W)+R(t)$. Their initial offset must be retained:
$b_j=\beta_j-R_{j,R}(W)$. Since $W\leq o$ and $B\leq o^2/1024$,
$|b_j|\leq o^2+oW+B\leq4o^2$. Only the $W-1$ internal gaps are used.
\end{proof}

\subsection{Forcing the palette inside the core}
Use
\[
T=2^{13}o^2,\qquad A_0=2^{13}(S+1)T,\qquad \theta_j=jT.
\]
For early--late edges let $u_{L,j}$ be the ordinary positive half-frame weight at early endpoint $j$, and $u_{R,j}$ the magnitude of the negative weight at its late partner. Choose
\begin{align}
C_{L,j}&=\ceil*{\frac{A_0+\theta_j-b_j}{u_{L,j}}},
& C_{L,j}u_{L,j}&=A_0+\theta_j-b_j+e_{L,j},\notag\\
C_{R,j}&=\ceil*{\frac{A_0+\theta_j}{u_{R,j}}},
& C_{R,j}u_{R,j}&=A_0+\theta_j+e_{R,j},
\label{eq:ordinary-counts}
\end{align}
where $0\leq e_{L,j},e_{R,j}<B$. Both targets are positive. These ceilings use only known integer weights and palette crossing counts.

\begin{lemma}[Palette forcing and selected residuals]\label{lem:palette}
Every minimizer defining $M_C$ agrees with $f^*$ on $\mathcal F$. Once the palette and fans are fixed, each selected edge has unary coefficient $r_e$ with $|r_e|<4B$; in particular $|r_e|\leq10B$.
\end{lemma}
\begin{proof}
For a nonselected early-$j$--late-$k$ edge, $j\ne k$, the ordinary coefficient is
\[
\theta_j-\theta_k-b_j+e_{L,j}-e_{R,k}.
\]
The first term has absolute value at least $T$. All adverse terms, including arbitrary remaining core colours, have magnitude at most
\begin{align*}
|b_j|+2B+2(4oW/B+4)B+\binom{o-2}{2}
&\leq\tfrac92o^2+8oW+10B\\
&<5o^2<T.
\end{align*}
The last strict bound uses $o\geq514W$ and $B\leq o^2/1024$. The rank term counts both incident endpoint families; nonincident programs contribute nothing. The binomial term bounds the full core crossing degree, without assuming any other palette colour. A wrong colour can therefore be improved by a single flip. The preferred colour is zero for $j>k$ and one for $j<k$.

For an edge on one side, the two ordinary fields have the same sign. Each ordinary multiplicity is at least $A_0/(2B)$ because its target is at least $A_0-4o^2$. Either half-frame then contributes magnitude at least $A_0/8$ from each parent. A side--dummy edge has one such contribution. All adverse rank and core terms have magnitude at most
\[
8oW+8B+\binom{o-2}{2}<o^2<A_0/8.
\]
Thus early-side edges prefer zero and late-side edges prefer one. This covers the palette simultaneously.

For the selected edge $p_{j,r}q_{j,t}$, the ordinary counts cancel $b_j$ in~\eqref{eq:corrected-field}, leaving exactly
\begin{equation}\label{eq:selected-residual}
r_e=\eps_{j,L}(r)+\eps_{j,R}(t)+e_{L,j}-e_{R,j}.
\end{equation}
All four errors have magnitude below $B$, proving the residual bound.
\end{proof}

This lemma is restricted to~\eqref{eq:core-energy}. Before fan freezing, a global optimum need not have this palette. The full-drawing lower bound will use only that its post-freezing core energy is at least $M_C$.

\subsection{Dummy edges and local constraints}
\begin{lemma}[Exact dummy elimination]\label{lem:dummy}
With the palette fixed, the optimized dummy-edge contribution, including its palette interactions, is
\begin{equation}\label{eq:dummy}
Z(D_m)+SW\binom{D_m}{3},
\end{equation}
independently of all selected and short-edge colours.
\end{lemma}
\begin{proof}
A dummy chord with ranks $a<b$ crosses no early--late selected chord and no endpoint-cluster short edge. Its palette neighbours are exactly the side--dummy edges whose dummy endpoint is strictly between $a,b$. It has $SW(b-a-1)$ such neighbours of each colour. Its palette cost is consequently the constant $SW(b-a-1)$ in either colour. There are no program-fan neighbours by Lemma~\ref{lem:nonincident}. Summing over $a<b$ gives $SW\binom{D_m}{3}$. What remains is the full convex drawing on $D_m$ points, with optimum $Z(D_m)$ by the two-page complete-graph theorem \cite{Abrego2013}. Its endpoint-alternation predicate agrees with that of convex straight-line chords.
\end{proof}

Let $Q_{\rm loc}$ be the crossing constraints internal to one selected bundle, between selected and short edges, and between two short edges. These are crossing-graph constraints, not removed geometric edges.

\begin{lemma}[Conservative local budget]\label{lem:local-count}
\begin{equation}\label{eq:local-budget}
|Q_{\rm loc}|\leq2SW^4.
\end{equation}
After deleting these constraints and fixing the palette, every short edge has only unary costs. The dummy subsystem remains independent, and the remaining selected-edge graph is that of Lemma~\ref{lem:source-amplification}.
\end{lemma}
\begin{proof}
In one selected bundle, an internal crossing is specified by two early and two late ranks, giving at most $W^4/4$ choices. For a source--short crossing, the short chord belongs to one of that source's two endpoint clusters and its endpoints surround the source endpoint. An upper bound, even without this surrounding condition, is
$2\binom W2 W^2\leq W^4$ per bundle. A short--short crossing must lie in one cluster; the two clusters contribute at most $2\binom W2^2\leq W^4/2$. The sum is at most $7W^4/4\leq2W^4$.

A short edge cannot cross a selected bundle avoiding its cluster, since the cluster is consecutive in cap order; short edges in distinct clusters cannot alternate. Thus the listed types exhaust its unfixed neighbours, except for dummy edges, which are noncrossing with it by Lemma~\ref{lem:dummy}. The remaining fixed neighbours give unary costs. Removing internal bundle crossings leaves exactly the source graph already identified.
\end{proof}

\subsection{A computable core baseline}
Delete $Q_{\rm loc}$ from~\eqref{eq:core-energy}. Fix the palette, set all selected colours to zero, independently minimize each short unary, and insert the dummy optimum~\eqref{eq:dummy}. Call this integer $C_0$, and define
\begin{equation}\label{eq:core-baseline}
K_C=C_0-\Delta|E(G)|+\sum_{e\in\mathcal S}\min\{0,r_e\}.
\end{equation}
All of these operations are explicit sums, crossing tests, and minima of two integers. In particular neither $g(G)$ nor $M_C$ is used.

For a fixed selected colouring $x$, let $\Psi(x)$ be the minimum of the deleted energy over the short and dummy variables, with the palette fixed. The exact identity is
\begin{align}
\Psi(x)
&=C_0-\Delta|E(G)|+\Phi_{\rm main}(x)+\sum_{e\in\mathcal S}r_ex_e\notag\\
&=K_C+\Phi_{\rm main}(x)
+\sum_{e\in\mathcal S}\bigl(r_ex_e-\min\{0,r_e\}\bigr).
\label{eq:core-identity}
\end{align}
Indeed, the short minima and dummy value are independent of $x$ after deletion. All remaining source-dependent fixed-neighbour costs have coefficient $r_e$. The all-zero selected colouring contributes exactly $\Delta|E(G)|$ to the main graph, explaining the single subtraction. No internal-bundle contribution is subtracted twice.

\begin{proposition}[Core encoding]\label{prop:core-encoding}
The baseline $K_C$ is computable in polynomial time, and
\begin{equation}\label{eq:core-bound}
K_C+\Delta g(G)\leq M_C\leq
K_C+\Delta g(G)+10BSW^2+2SW^4.
\end{equation}
\end{proposition}
\begin{proof}
Every core minimizer has the palette by Lemma~\ref{lem:palette}. Deleting nonnegative local terms and applying~\eqref{eq:core-identity} gives the lower bound, since $\Phi_{\rm main}(x)\geq\Delta g(G)$ and each residual excess is nonnegative.

For the upper bound take a class-uniform selected colouring attaining $\Delta g(G)$, independently minimize every short unary, and choose an optimal dummy colouring. These choices coexist in one complete core colouring. Each selected unary exceeds its separate minimum by at most $|r_e|\leq10B$, and $|\mathcal S|=SW^2$. Restoring all local constraints costs at most $2SW^4$ by Lemma~\ref{lem:local-count}. The formula~\eqref{eq:core-baseline} proves computability with polynomially many exact operations. The optimizers just used certify the inequalities and are not computed by the reduction.
\end{proof}

The bound $10BSW^2$ leaves some slack. The exact residual~\eqref{eq:selected-residual} permits $4BSW^2$ in this unordered-edge presentation. The formal proof uses $8SH^{13}\leq10SH^{13}$, which is the same stated budget because $B=H^5$ and $W=H^4$.

\begin{proposition}[Program count]\label{prop:program-count}
The number of required program points satisfies
\begin{equation}\label{eq:program-count}
s\leq\frac{16SA_0}{B}+\frac{8SoW}{B}+10S
\leq2^{80}(S+1)^4H^3.
\end{equation}
In particular $s\leq o$. All multiplicities are computable before final general-position insertion.
\end{proposition}
\begin{proof}
Each target in~\eqref{eq:ordinary-counts} is at most $A_0+ST+4o^2<2A_0$. A weight is at least $B/4$, so each endpoint uses at most $8A_0/B+1$ ordinary programs. Lemma~\ref{lem:rank-correction} uses at most $4oW/B+4$ rank programs per endpoint. Multiplication by $2S$ gives the first bound.

For $d=S+1$, $o=514SH^4<2^{10}dH^4$ and
$A_0=2^{26}do^2<2^{46}d^3H^8$.
Thus the first two terms are below $2^{50}d^4H^3$ and $2^{13}d^2H^3$, respectively, and $10S$ is smaller still. Their sum is below $2^{80}d^4H^3$. Since $H=2^{20000}d^{40}$, this is smaller than $514SH^4=o$ for $S\geq1$. Each step uses only the individually known response weights; placing the resulting number of copies in their open neighbourhoods preserves those weights.
\end{proof}

\section{The complete drawing and the decision reduction}\label{sec:reduction}

All coordinates and crossing predicates in this section refer to the final general-position drawing from Proposition~\ref{prop:realization}. Its total number of points is $N=B+o+s$; the comparison grid has $N_g=Bo$ vertices. The notation $\Phi_{AB}$ means the sum of monochromatic constraints with one endpoint in each of the distinct crossing-graph vertex classes $A,B$; $\Phi_A$ counts unordered pairs within $A$.

\subsection{A common computable background}
Let $G_0$ be the crossing graph with $Y$ deleted. Its remaining constraints are precisely
\[
X\!:\!X,\quad X\!:\!T_f,\quad T_f\!:\!T_f,\quad C\!:\!C,\quad C\!:\!T_f,\quad I\!:\!I.
\]
Write $c^+$ for the canonical grid colouring and $t^+$ for the independently computed program-fan responses, with colour zero at a zero field. Define
\begin{align}
E_X(x;t^+)&=\Phi_X(x)+\Phi_{XT_f}(x,t^+),\notag\\
d_v&=E_X(c^+\oplus\{v\};t^+)-E_X(c^+;t^+),\notag\\
L_+&=E_X(c^+;t^+)+\sum_{v\in X}\min\{0,d_v\},\notag\\
K_0&=Z(B)+\Phi_{T_f}(t^+).
\label{eq:background}
\end{align}
The fan--fan term is evaluated after insertion, including every crossing that changed during program placement. These constants require no optimal grid or fan colouring.

\begin{proposition}[Decoupling with the grid fixed]\label{prop:decoupling}
Suppose the hypotheses of Theorem~\ref{thm:grid} and Proposition~\ref{prop:freezing} hold with grid dimensions $B,o$, grid size $N_g=Bo$,
\[
\rho=1/B+1/o,\quad a=\alpha+\rho\leq2^{-4096},\quad
Q=12s/o,\quad b=104a+\beta\leq1/100.
\]
Set $E_g=2^{32}N_g^2a^3$ and $E_{\rm fr}=2Q(b+\rho)bN_g^2$. Then
\begin{align}
K_0+L_++M_C-E_{\rm fr}-E_g
&\leq\OPT(G_0)\leq K_0+L_++M_C+E_g,\label{eq:decoupling-zero}\\
K_0+L_++M_C-E_{\rm fr}-E_g
&\leq\OPT(P)\leq K_0+L_++M_C+E_g+R_Y.
\label{eq:decoupling-full}
\end{align}
\end{proposition}
\begin{proof}
Take an actual optimum of $G_0$. Conditional on every other colour, its grid minimizes a unary-perturbed comparison energy with unary differences bounded by $\alpha N_g$. Theorem~\ref{thm:grid} supplies a canonical phase and a localized flip set of size at most $52aN_g$. If necessary complement every colour in the whole graph to make the phase positive. There are no prescribed colours, and complementation preserves every monochromatic count.

Keep this grid unchanged. By Lemma~\ref{lem:soft-fields}, the probe soft-count hypothesis holds. Proposition~\ref{prop:freezing} changes only the program fans to $t^+$, with cost increase at most $E_{\rm fr}$. The grid retains its localized flip set, so the localized lower estimate in Theorem~\ref{thm:grid}, now for the frozen unary energy in~\eqref{eq:background}, gives at least $L_+-E_g$. The current core colouring costs at least $M_C$, irrespective of whether its palette was forced before the fans were frozen. The frame class costs at least $Z(B)$, and the fan--fan cost is exactly $\Phi_{T_f}(t^+)$. These terms exhaust $G_0$. Thus the altered colouring has cost at least $K_0+L_++M_C-E_g$ but at most $\OPT(G_0)+E_{\rm fr}$, proving the lower bound.

For the upper bound, use the same $t^+$, the grid colouring
$c^+\oplus\{v:d_v<0\}$, a core minimizer, and a convex frame optimum. The prescribed-phase upper witness in Theorem~\ref{thm:grid} has grid cost at most $L_++E_g$. All these choices coexist, since the edge classes are disjoint and no further cross-class constraints remain. This proves~\eqref{eq:decoupling-zero}. Restoring $Y$ with~\eqref{eq:nonfan-delete} proves~\eqref{eq:decoupling-full}. Optimizers are used only as witnesses for inequalities.
\end{proof}

\subsection{One parameter choice verifies every hypothesis}
Use the parameters of~\eqref{eq:parameters}, namely
\[
q=2^{20}(n+1),\quad S=nq,\quad d=S+1,\quad
H=2^{20000}d^{40},\quad W=H^4,\quad B=H^5,\quad o=514SW.
\]
Take
\begin{equation}\label{eq:field-scales}
\alpha=\beta=\frac{2^{84}d^4}{H},\qquad
a=\alpha+\rho\leq\frac{2^{85}d^4}{H},\qquad
b=104a+\beta.
\end{equation}
Proposition~\ref{prop:program-count} gives $s\leq o$ and the convenient bounds
\begin{equation}\label{eq:all-sizes}
N_g\leq2^{10}dH^9,\qquad p=Bs\leq2^{80}d^4H^8,\qquad
c\leq2^{19}d^2H^8,\qquad h\leq2os\leq2^{91}d^5H^7.
\end{equation}
Also $B\geq2o$, $B\leq o^2/1024$, $o\geq1000$, $8\mid B$, and $W/o\leq1/514$. These verify the geometric and core hypotheses. Each follows directly by substituting $W=H^4$ and $B=H^5$ and using $H=2^{20000}d^{40}$, $d\geq1$, $S\geq1$.

In $G_0$, a grid vertex has at most $p$ outside neighbors: it has no neighbors in $C$ or $I$. A program fan has at most $p+c$ neighbors outside $X$, since $I$ is isolated. Thus
\[
p/N_g=s/o\leq\alpha,\qquad
(p+c)/N_g\leq\beta,\qquad
Q=12s/o\leq\alpha\leq a.
\]
For the second inequality, $c/N_g<o/(2B)$ and the spare factor in $2^{84}$ dominates this term together with $s/o$. The mesh satisfies $\rho\leq\alpha$. Finally the scale of $H$ gives $a<2^{-4096}$ and $b<1/100$. The two phases, the localized grid estimate, and fan freezing therefore apply to this one actual drawing.

Substitution in their error bounds gives
\begin{equation}\label{eq:individual-errors}
E_g\leq2^{307}d^{14}H^{15},\qquad
E_{\rm fr}\leq2^{290}d^{14}H^{15},\qquad
R_Y\leq2^{182}d^{10}H^{15}.
\end{equation}
For the freezing term, $Q\leq a$, $b\leq105a$, and $b+\rho\leq106a$, whence
$E_{\rm fr}\leq2\cdot105\cdot106N_g^2a^3<2^{15}N_g^2a^3$.
For $R_Y$, use $p+c\leq2^{81}d^4H^8$ and~\eqref{eq:all-sizes}. The resulting product is at most $2^{172}d^9H^{15}$, and $h^2/2\leq2^{181}d^{10}H^{14}$, both covered by the stated bound. Consequently
\begin{equation}\label{eq:algebraic-error}
E_{\rm fr}+2E_g+R_Y<2^{310}d^{14}H^{15}.
\end{equation}

\subsection{Quantitative reduction and threshold}
\begin{theorem}[A computable source-value window]\label{thm:quantitative-reduction}
Given a represented permutation graph $G$ with $n\geq1$, the construction computes a general-position rational point set $P$, a baseline $K_{\rm base}$, and a positive integer $\Delta=q^2H^{16}$, with explicit nonnegative errors such that
\begin{equation}\label{eq:final-window}
K_{\rm base}+\Delta g(G)-E_-
\leq\OPT(P)\leq
K_{\rm base}+\Delta g(G)+E_+,\qquad
2^{18}(E_-+E_+)<\Delta.
\end{equation}
All output data and threshold constants have polynomial binary length and are computable in deterministic polynomial time. The points can instead be returned with integer coordinates.
\end{theorem}
\begin{proof}
Put
\begin{align}
K_{\rm base}&=K_0+L_++K_C,\notag\\
E_-&=E_{\rm fr}+E_g,\notag\\
E_+&=E_g+R_Y+10BSW^2+2SW^4.
\label{eq:final-errors}
\end{align}
The core bound~\eqref{eq:core-bound} and Proposition~\ref{prop:decoupling} give the two-sided interval.

Divide its total error by $\Delta=q^2H^{16}$. The part in~\eqref{eq:algebraic-error} is less than
\[
\frac{2^{310}d^{14}}{q^2H}
\leq2^{-19690}d^{-26}<2^{-40}.
\]
The selected unary term satisfies
\[
\frac{10BSW^2}{\Delta}=\frac{10n}{qH^3}<2^{-40},
\]
and the conservative local term satisfies
\[
\frac{2SW^4}{\Delta}=\frac{2n}{q}<2^{-19}.
\]
Hence the total is less than $2^{-39}+2^{-19}<2^{-18}$, proving the required strict window. No sharper local enumeration or narrower aggregate window is used.

To compute the baseline, first enumerate proper crossings on the actual final coordinates by exact determinants. Evaluate $c^+$, each fan response, the actual fan--fan cost, and each single-flip difference $d_v$. Evaluate $K_C$ by~\eqref{eq:core-baseline}: set the palette, use the all-zero source colour as a reference, minimize each short unary separately, insert the known dummy value, subtract the explicit $\Delta|E(G)|$, and add the residual unary minima. Finally insert $Z(B)$. These are finite sums and minima of two known integers. The source optimum, core optimum, and target optimum never occur as computational subroutines.

The output point count is $N=B+o+s\leq3H^5$. Since $S=2^{20}n(n+1)$, it is bounded by a fixed polynomial of degree $400$ in $n$. Proposition~\ref{prop:realization} proves polynomial coordinate bit length and construction time; it also gives the integer-coordinate version by a common positive scaling.
\end{proof}

We make the bit bound for the decision threshold explicit as well. For actual total point count $N$, all grid, fan, and core edge counts are at most $N^2$. Each unnormalized crossing contribution is zero or one, and all fields are differences of finite such sums. The baseline formula~\eqref{eq:background} and~\eqref{eq:core-baseline} therefore gives the deliberately loose bound
\[
|K_{\rm base}|\leq16N^6.
\]
One may also see this by computing each baseline from its fixed colouring cost and at most $N^2$ single-edge unary corrections, each bounded by a fixed multiple of $N^2$; $\Delta|E(G)|$ is the all-zero main crossing count, at most $N^4$. Since $q,H\leq N$, we have $\Delta\leq N^{18}$, and $|E(G)|\leq N^2$.

For a source MaxCut threshold $k\in\Z$, put $t=|E(G)|-k$ and output
\begin{equation}\label{eq:decision-threshold}
\kappa=\floor*{K_{\rm base}+\Delta t+\Delta/4}.
\end{equation}
In particular
\begin{equation}\label{eq:threshold-bits}
|\kappa|\leq32N^{20}(|k|+1),
\end{equation}
so the threshold bit length is $O(\log N+\operatorname{bits}(k))$. The baseline is integral in the counting presentation above; the floor also accommodates its equivalent rational presentation in the formalization.

\begin{proof}[Proof of Theorem~\ref{thm:main}]
Use the NP-hard represented permutation-graph problem with threshold $k$, or equivalently $g(G)\leq t$. For $n\geq1$ apply Theorem~\ref{thm:quantitative-reduction} and output~\eqref{eq:decision-threshold}. If $g(G)\leq t$, then
\[
\OPT(P)\leq K_{\rm base}+\Delta t+E_+
<K_{\rm base}+\Delta t+\Delta/4.
\]
Since the objective is an integer, it is at most $\kappa$. If $g(G)>t$, integrality gives $g(G)\geq t+1$, and
\[
\OPT(P)\geq K_{\rm base}+\Delta(t+1)-E_-
>K_{\rm base}+\Delta t+\Delta/4\geq\kappa.
\]
Thus the same point set and threshold give both directions.

Empty source inputs can be decided directly. A fixed yes instance consists of two distinct integer points and threshold zero. A fixed no instance is the five points $(i,i^2)$, $0\leq i<5$, with threshold zero: their diagonal crossing graph is a five-cycle and cannot be properly two-coloured. If the target convention permits only nonnegative thresholds and the computed $\kappa$ is negative, output this fixed no instance. A source yes instance would otherwise imply the impossible inequality $0\leq\OPT(P)\leq\kappa<0$. Thresholds outside the source edge-count range may also be handled directly, though the displayed equivalence already covers them.

Proposition~\ref{prop:realization} and~\eqref{eq:threshold-bits} prove a deterministic polynomial-time many-one reduction with a single complete geometric target instance. Every target edge colour is a free certificate bit. The palette and frozen colours enter the proof through minimizer properties and one compatible upper-bound colouring, not through precolouring restrictions on the output. Scaling produces the same hardness on integer coordinates. Membership in NP follows by assigning colours to all $\binom N2$ edges and testing the $O(N^4)$ edge pairs with exact orientation determinants. This proves NP-completeness and, consequently, NP-hardness of exact minimization.
\end{proof}

\begin{corollary}
If collinear input points are permitted while the objective still counts only proper crossings using strict orientation tests, the decision problem remains NP-complete.
\end{corollary}
\begin{proof}
The general-position hard instances are also inputs in the larger class. The same polynomial verifier excludes every zero orientation determinant and therefore handles the stipulated crossing convention on degenerate inputs. No perturbation of an arbitrary degenerate input is used.
\end{proof}

\section{Formal verification and reproducible computations}\label{sec:formalization}

The accompanying Lean development formalizes the new mathematical
construction and its final interface with the complete geometric graph
problem. The principal declaration is
\begin{center}
\texttt{Q4.Candidate.ConstructedReduction.}\\
\texttt{constructed\_maxcut\_instance}.
\end{center}
For every $n\geq1$, permutation $\pi$ and integer source threshold $k$,
it produces one point count $N$, one coordinate function
$\mathrm{coords}:\mathrm{Fin}(N)\to\Q^2$, and one integer threshold
$\kappa$. The same witness satisfies injectivity, general position,
coordinate encoding bounds and the equivalence
\[
 \exists\chi\;M(\mathrm{coords},\chi)\leq\kappa
 \quad\Longleftrightarrow\quad
 \exists\text{ a cut of the inversion graph of }\pi\text{ of size at least }k.
\]
The colours on the left range over every edge of $K_N$ freely.
The formal construction gives, with
$\widehat H=2^{20840}(n+1)^{80}$,
\[
 N\leq3\widehat H^5,
 \qquad L_{\mathrm{coord}}\leq2^{17}\widehat H^{10},
\]
where each coordinate has an integer numerator and nonzero denominator
of absolute value at most $2^{L_{\mathrm{coord}}}$.

\subsection{What the theorem assumes and what it proves}

The final Lean declaration has the explicit mathematical hypothesis
\begin{center}
\texttt{forall m, ConvexOptimum m (knownZ m)}.
\end{center}
Here \texttt{ConvexOptimum} asserts both a lower bound for every free
two-colouring of the chords and existence of an attaining colouring;
\texttt{knownZ} is the value $Z$ above. This hypothesis is supplied in the
mathematical proof by the established two-page theorem
\cite{Abrego2013}. Its conversion to the actual frame and dummy chord
objectives is formalized. The literature theorem itself is not reproved
in this development and is not introduced as a custom Lean axiom.

All other construction obligations in the final declaration are
instantiated: actual program multiplicities, geometry, general-position
placement, palette forcing in the appropriate core optimization,
selected-edge residual fields, source amplification, the complete edge
partition, the computable background, and the strict error gap.
In particular, an unknown source optimum or background optimum is not
an input to the reduction.

\begin{center}
\small
\begin{tabular}{p{.29\linewidth}p{.64\linewidth}}
\hline
Proof component & Principal module families\\
\hline
Analytic and discrete stability &
\texttt{TwoPhaseStability}, \texttt{GridStepBridge},
\texttt{GridLocalization}\\
Complete energy and freezing &
\texttt{CompleteProbeFreezing}, \texttt{CompleteClassEnergy},
\texttt{CompleteBackground}\\
Actual geometric programs &
\texttt{ProgramGeometry}, \texttt{ActualProgramGeometry},
\texttt{ActualRecipe}, \texttt{ProgramPlacement}\\
Free core and source signal &
\texttt{ActualCatalogBounds}, \texttt{ActualSelectedUnary},
\texttt{MainSourceEncoding}, \texttt{CoreValueBridge}\\
Effective baseline and encoding &
\texttt{ComputedBackground}, \texttt{ComputedCore},
\texttt{RationalEncoding}, \texttt{CandidateSize}\\
Final complete-input equivalence &
\texttt{CompleteInput}, \texttt{ActualReductionValue},
\texttt{SourceCutBridge}, \texttt{ConstructedReduction}\\
\hline
\end{tabular}
\end{center}

The development uses the conservative local bound $2SW^4$ and proves
$2^{18}(E_-+E_+)<\Delta$. These are the bounds used in the final
reduction here. No stronger constant from an earlier draft is needed.

The source NP-hardness theorem \cite{Figueiredo2023}, complexity classes,
machine running times, membership in NP, and the separate output-threshold
bit bound are not formalized in Lean. Their mathematical proofs and
references are given in this paper and in the effective-reduction notes.
In particular, a small-coordinate existence theorem alone would not
prove the required polynomial construction time; the explicit finite
selection algorithm supplies this additional fact.

\subsection{Verification procedure}

The public proof artifact contains exactly the transitive project import
closure of the final theorem: 123 construction modules and four shared
foundation modules, with a separate explicit aggregate entry point.
The Lean version is fixed to 4.34.1 and the mathlib dependency and its
transitive revisions are fixed by the Lake manifest.

The verification script first compares the source set with this closure
and its SHA256 manifest. It checks for proof placeholders and custom
axioms, builds the aggregate entry point, elaborates every source file
individually, and checks the outputs of the retained explicit
\texttt{\#print axioms} commands, including the final theorem. It also
checks that source and configuration hashes did not change during the
run. The permitted kernel axiom dependencies are
\texttt{propext}, \texttt{Classical.choice}, and \texttt{Quot.sound}.
The axiom queries cover the designated key declarations; they are not
described as a separate enumeration of every declaration in mathlib.
The trust boundary includes Lean, the pinned dependencies and their
compiled artifacts, and the ordinary compiler and hardware stack.

\subsection{Finite diagnostics and review scope}

Four standard-library Python programs provide exact finite diagnostics:
the six-point completion obstruction in Appendix~\ref{app:completion};
core combinatorics for all six orders of three source vertices; two
reduced geometric realizations with actual general-position insertion
and crossing-table checks; and rational or integer checks of selected
parameter inequalities. The programs use exact arithmetic. Their
recorded inputs and expected outputs are distributed with the source.
They do not generate or solve the extremely large target instances of
the reduction and are not used as proofs of statements for unbounded
input sizes.

AI systems assisted the exploration, formalization and preparation of
the manuscript. Independent LLM-assisted reviews were used to examine
definitions, interfaces and the effective construction, followed by
kernel checks of the formal development. The artifact records this
scope explicitly; these reviews do not constitute human journal peer
review. The repository retains the final proof and useful diagnostics,
with intermediate exploratory records kept outside the public source
set.

\appendix
\section{A direct-completion obstruction}\label{app:completion}

The following example explains why the contribution of added edges must
be controlled in a completeness-preserving reduction.

\begin{proposition}
Completing a segment family to the complete geometric graph on its
endpoints can produce an optimum whose restriction to the original
family is not optimal. The minimum contribution of the added edges can
also depend on the prescribed source colours.
\end{proposition}
\begin{proof}
Take $p_i=(i,i^2)$ for $0\leq i<6$ and abbreviate $p_ip_j$ by $ij$.
The source edges are $01,24,35$. Their only crossing pair is $24,35$,
so the source optimum is zero.

Every five-point subset is strictly convex. The crossing graph of its
five diagonals is an odd cycle, so every two-colouring contains a
monochromatic crossing. There are six five-point subsets, and every
crossing pair belongs to exactly two of them. Therefore the complete
six-point graph has at least three monochromatic crossing pairs in
every colouring.

Colour $03,04,13,14$ blue and all remaining edges red. Its monochromatic
crossing pairs are exactly $(02,15),(03,14),(24,35)$, so it is optimal
for the complete graph. Its restriction to the source has cost one.
In particular, when $24$ and $35$ have the same colour, the minimum
additional cost is two.

Alternatively colour $14,15,24,25$ blue and the other edges red.
The monochromatic crossing pairs are $(02,13),(04,35),(14,25)$.
This is again a complete optimum, but $24$ and $35$ now have different
colours. The source cost is zero and the minimum additional cost is
three. Hence the extra optimum is not independent of the source
colouring.
\end{proof}

\bibliographystyle{plain}
\bibliography{references}
\end{document}
